\documentclass[11pt,oneside]{report}
% Shared preamble for the HKJC win-model primer.
\usepackage[utf8]{inputenc}
\usepackage[T1]{fontenc}
\usepackage{lmodern}
\usepackage[a4paper,margin=1in]{geometry}
\usepackage{amsmath,amssymb,amsthm,mathtools}
\usepackage{booktabs}
\usepackage{array}
\usepackage{longtable}
\usepackage{enumitem}
\usepackage{xcolor}
\usepackage{microtype}
\usepackage{siunitx}
\usepackage{graphicx}
\usepackage[hidelinks]{hyperref}
\usepackage{fancyhdr}

\hypersetup{colorlinks=true,linkcolor=blue!45!black,citecolor=blue!45!black,urlcolor=blue!45!black}

% ---- notation shortcuts -------------------------------------------------------
\newcommand{\vc}{{v_{c}}}
\newcommand{\Wp}{{W^{\prime}}}
\newcommand{\Tfin}{T}
\newcommand{\Leff}{{L_{\mathrm{eff}}}}
\newcommand{\Ddiff}{D}
\newcommand{\argmin}{\operatorname*{arg\,min}}
\newcommand{\IG}{\mathcal{IG}}
\newcommand{\Norm}{\mathcal{N}}
\newcommand{\E}{\mathbb{E}}
\newcommand{\Var}{\operatorname{Var}}
\newcommand{\Cov}{\operatorname{Cov}}
\newcommand{\given}{\,\vert\,}
\newcommand{\real}{\mathbb{R}}
\newcommand{\pcterm}[1]{\textcolor{blue!45!black}{\textbf{#1}}}
\newcommand{\code}[1]{\texttt{#1}}

\theoremstyle{definition}
\newtheorem{principle}{Principle}
\newtheorem{remark}{Remark}[chapter]
\newtheorem{definition}{Definition}[chapter]
\newtheorem{proposition}{Proposition}[chapter]


\title{\Huge\bfseries A Physics-Grounded Model of\\[2mm] Hong Kong Horse-Racing Win Probability\\[6mm]
  \Large A First-Principles Treatment of the Generative Law,\\ the Latent Energetics, and the Calibration Programme}
\author{The \texttt{hkjc-tools} modelling stack\\[2mm]\normalsize Internal technical primer}
\date{June 2026}

\begin{document}

\pagenumbering{roman}
\maketitle

\begin{abstract}
\noindent This primer is a complete, self-contained account of the production win-probability
model for Hong Kong Jockey Club (HKJC) thoroughbred racing as it stands in mid-2026. The model
is not a regression on engineered covariates; it is a \emph{generative} order-statistic model in
which each horse's finishing time is the first-passage time of a drift--diffusion process, the
winner is the runner whose first passage is earliest, and every input couples to one of three
mechanistically-interpretable parameters of that process: the aerobic \emph{critical velocity}
$\vc$, the finite \emph{anaerobic reserve} $\Wp$, and the \emph{diffusion} $\Ddiff$ governing
race-to-race variance. We develop the model from the generative spine downward: the
Inverse-Gaussian (Wald) first-passage law and its closed-form Plackett--Luce surrogate
(Chapter~\ref{ch:gen}); the inverse problem of recovering $(\vc,\Wp)$ from sectional split times
(Chapter~\ref{ch:energetics}); the hierarchical/empirical-Bayes population that turns noisy
per-run estimates into leak-free, partially-pooled, rating-conditioned per-horse abilities
(Chapter~\ref{ch:pop}); the full leak-free feature layer --- recency-weighted form, jockey and
trainer signals, running-style and positional traits, velocity speed-figures, barrier-trial
debut form, trajectory-shape traits, and rating-trajectory process features
(Chapter~\ref{ch:features}); the coherent precision-fusion that folds heterogeneous velocity
observations into a single $\vc$ posterior at each channel's own precision
(Chapter~\ref{ch:fusion}); the relational opponent-adjusted Massey ``collateral-form'' rating
that operationalises the competition axis (Chapter~\ref{ch:competition}); the geometry of the
movable running rail and the bend ground-loss that draw imposes on both the drift and the
diffusion (Chapter~\ref{ch:geometry}); and the assembled drift/reserve and the Monte-Carlo
predict loop (Chapter~\ref{ch:assembly}). Chapter~\ref{ch:theory} sets out the theoretical
frontier: the race as a mean-field differential game (the longitudinal pace game, the lateral
geometry--control game, and the temporal phase structure of the trip). Chapter~\ref{ch:method}
states the calibration philosophy and the paired fixed-seed conversion test that gates every
change. Chapter~\ref{ch:state} reports the current calibration, the empirical results, and the
open research queue. The throughline is a single discipline: \emph{fit the market} as a
goodness-of-fit objective, regularised by physics and game theory so that every degree of
freedom is a named mechanism rather than a free parameter.
\end{abstract}

\tableofcontents

\cleardoublepage
\pagenumbering{arabic}

\chapter{Introduction: the object, the objective, and the discipline}
\label{ch:intro}

\section{What problem is being solved}

A horse race is, formally, the realisation of an \emph{order statistic}. For a field of $N$
runners, each horse $i$ produces a finishing time $\Tfin_i\in\real_{>0}$, and the result of the
race --- the only thing that pays a win bet --- is the identity of the minimiser,
\begin{equation}
  \text{winner} \;=\; \argmin_{i\in\{1,\dots,N\}} \Tfin_i .
  \label{eq:argmin}
\end{equation}
Everything in this document is, ultimately, a theory of the joint law of
$(\Tfin_1,\dots,\Tfin_N)$ and of the marginal it induces on the winner through
\eqref{eq:argmin}. We do not model ``the probability that horse $i$ wins'' directly with a
classifier; we model the \emph{generative process of finishing times} and read the win
probability off it as
\begin{equation}
  p_i \;=\; \Pr\!\left[\Tfin_i < \Tfin_j \ \forall j\neq i\right].
  \label{eq:winprob}
\end{equation}
This is a deliberate and consequential choice. It means that the same machinery yields the full
finishing-order distribution (hence place and exotic prices), that distributional facts about
finishing times --- in particular \emph{variance}, not only the mean --- are first-class
citizens of the model, and that every covariate must earn its place by changing the law of a
time, not by correlating with an outcome.

\section{The objective: fit the market as goodness-of-fit}

The betting market is the most efficient publicly-available aggregator of information about
\eqref{eq:winprob}. We therefore take the market's implied (de-vigged) win probabilities as the
\emph{yardstick} and pose the modelling problem as one of \emph{goodness-of-fit to the market},
measured strictly out-of-sample and walk-forward. The headline metric is the fraction of the
market's log-likelihood advantage over a uniform baseline that the model recovers,
\begin{equation}
  \%\text{ of market} \;=\;
   \frac{\mathrm{LL}_{\text{uniform}} - \mathrm{LL}_{\text{model}}}
        {\mathrm{LL}_{\text{uniform}} - \mathrm{LL}_{\text{market}}},
  \label{eq:pctmarket}
\end{equation}
where each $\mathrm{LL}$ is the mean log-probability assigned to the actual winners over a held-out
set of races, $\mathrm{LL}_{\text{uniform}}$ uses $p_i = 1/N$, and $\mathrm{LL}_{\text{market}}$
uses the de-vigged closing line. A score of $0\%$ is no better than ignorance; $100\%$ is parity
with the closing market. At project inception the model sat near $48\%$; at the time of writing it
is at approximately $65\%$ at Sha Tin (ST) and $60\%$ at Happy Valley (HV).

Three points about this objective deserve emphasis, because they shape every later decision.

\begin{principle}[Goodness-of-fit, not edge]
$\%$ of market is a \emph{goodness-of-fit} statistic. Closing the gap to the market \emph{is}
the modelling job. The model is not a search for mispricings; whether any residual gap to the
market is tradable is a separate, downstream, user-owned decision. The modelling work is pure
distributional fit.
\end{principle}

\begin{principle}[There is always a next mechanism]
A residual between model and market is treated as a \emph{missing mechanism} until proven
otherwise. The bar for declaring a residual to be genuinely unmodellable ``private information''
is high, because that same reasoning would have terminated the programme near $10\%$ skill.
Understanding \emph{what the market knows that the model does not} is the engine of progress, not
an excuse to stop.
\end{principle}

\begin{principle}[Calibration is part of the objective]
Fit is judged jointly on discrimination (\,$\%$ of market) \emph{and} calibration (the expected
calibration error, ECE). A change that sharpens discrimination while degrading calibration is not
an improvement. The current model's ECE is $\approx 0.006$, i.e.\ predicted win frequencies match
realised ones to well under a percentage point on average.
\end{principle}

\section{The discipline: first-principle regularisation}

The defining methodological commitment of this project is that the model is regularised
\emph{by physics and game theory, not by feature count or by a penalty norm}. Every input must
couple to the generative law through a named, plausible mechanism --- a term in the equation of
motion of a horse (the ``plant''), or a parameter of the jockey/trainer control policy. An
arbitrary regressor that merely correlates with winning is forbidden, however predictive it looks
in-sample, because it both overfits and breaks the chain of physical reasoning that has carried
the model this far.

This has a sharp operational form, developed fully in Chapter~\ref{ch:assembly}: the
drift--diffusion law exposes exactly \emph{three independent knobs} a feature may load:
\begin{enumerate}[label=(\roman*)]
  \item the aerobic ceiling $\vc$ (how fast the horse can sustainably go),
  \item the anaerobic reserve $\Wp$ (the finite ``kick'' it can spend above $\vc$), and
  \item the diffusion $\Ddiff$ (the race-to-race consistency / variance of its time).
\end{enumerate}
A feature is correctly \emph{structured} only when it loads the knob its mechanism actually acts
on. Using the right \emph{signal} is necessary but not sufficient; it must enter in the right
\emph{place}. A recurring lesson --- documented repeatedly in what follows --- is that a feature
which fails to convert under one structural coupling is not a dead mechanism but a wrongly-placed
one: scalar opponent-strength (Elo) failed, but the same competition axis expressed
\emph{relationally} (Chapter~\ref{ch:competition}) converts strongly; draw failed as a mean
effect on $\vc$ but converts as a variance effect on $\Ddiff$ (Chapter~\ref{ch:geometry}).

\begin{principle}[Theory leads, statistics gate]
Theory chooses \emph{which} structure to try; statistics only gate \emph{whether that particular
structure} converts out-of-sample. A dead first attempt is not a dead axis. Conversely, a
parked verdict can go stale when the baseline or the generative law beneath it changes, and must
be re-checked against the current production path before it is trusted.
\end{principle}

\section{Leak-freeness as a non-negotiable invariant}

A backtest number is worthless --- worse, actively misleading --- if any feature has seen the
future. The entire feature layer is therefore built to a single invariant: every per-horse
quantity used to predict a race on date $d$ is a function of that horse's \emph{strictly prior}
runs only. Concretely this is enforced two ways, both of which recur throughout
Chapters~\ref{ch:pop}--\ref{ch:competition}:
\begin{itemize}
  \item \textbf{Expanding means shifted by one.} A trailing trait is the cumulative sum over a
        horse's runs \emph{minus the current run}, divided by the count \emph{minus one}:
        $\bar{x}_{<t} = (\sum_{s\le t} x_s - x_t)/(n_{\le t}-1)$. The current observation never
        enters its own feature value.
  \item \textbf{Strict as-of joins.} Where a feature lives on a separate timeline (barrier
        trials, rating updates), it is joined backward with \code{allow\_exact\_matches=False},
        so a race on date $d$ sees only updates strictly before $d$.
\end{itemize}
Contemporaneous \emph{boundary conditions} that are genuinely known before the race --- the draw,
the track geometry and rail position, the field composition, the declared carried weight --- are
public pre-race information and are allowed. The discipline is enforced mechanically by an
\code{AsOfView} wrapper and an \code{assert\_no\_lookahead} check that is run on the fit/predict
pipeline before any backtest number is trusted.

\section{How to read this document}

The chapters are ordered from the generative spine outward. Chapter~\ref{ch:gen} writes down the
finishing-time law and its surrogate. Chapter~\ref{ch:energetics} recovers the latent energetics
from data. Chapter~\ref{ch:pop} builds the partially-pooled ability population. Chapters
\ref{ch:features}--\ref{ch:competition} are the feature layer: the channels that nudge a horse's
ability away from its population prior, and the principled way to combine them.
Chapter~\ref{ch:geometry} adds the arena. Chapter~\ref{ch:assembly} assembles everything into the
production drift and the predict loop. Chapter~\ref{ch:theory} is the theory of what is not yet
modelled --- the interaction games --- and is the source of the forward research queue.
Chapters~\ref{ch:method} and~\ref{ch:state} are the method and the current state. Appendices
collect the constants table and the data surface.

Throughout, the reader should keep the three knobs and the order statistic in mind: they are the
two ideas from which everything else is derived.

\chapter{The generative law: first-passage finishing times}
\label{ch:gen}
\label{ch:generative}

\section{From a velocity profile to a finishing time}

Fix a race of geometric length $L$ (metres). A horse covers the trip by accumulating distance at
a time-varying velocity $v_i(t)$, so its finishing time is the first time its accumulated
distance reaches $L$:
\begin{equation}
  \Tfin_i \;=\; \inf\Big\{\, t \ge 0 : \textstyle\int_0^t v_i(s)\,ds \;\ge\; L \,\Big\}.
  \label{eq:firstpassage-def}
\end{equation}
This is already a first-passage statement: the finishing time is the hitting time of a level $L$
by the integrated-velocity process. Everything in the theory of racing is, in the end, a theory
of \emph{what sets $v_i(t)$} --- how much speed the engine can produce, how that capacity is
deployed across the trip, how the arena and the rival field perturb it, and what is irreducibly
random. The chapters that follow populate $v_i(t)$ with mechanism. The present chapter fixes the
\emph{stochastic} structure of \eqref{eq:firstpassage-def} that turns a per-horse velocity model
into win probabilities.

\section{The drift--diffusion idealisation and the Wald law}

The production model idealises the accumulated-distance process $X_i(t) = \int_0^t v_i\,ds$ as a
Brownian motion with drift: in increments,
\begin{equation}
  dX_i \;=\; v_i\, dt \;+\; \sqrt{2\Ddiff_i}\; dW_i,
  \qquad X_i(0)=0,
  \label{eq:sde}
\end{equation}
where $v_i>0$ is the (effective, trip-averaged) drift velocity, $\Ddiff_i>0$ is a per-horse
diffusion coefficient with units of $\mathrm{m^2/s}$, and $W_i$ is a standard Wiener process. The
finishing time is the first passage of $X_i$ to the absorbing level $L$,
$\Tfin_i = \inf\{t : X_i(t)=L\}$. It is a classical result that the first-passage time of a
drift--diffusion to a fixed positive level follows the \emph{Inverse-Gaussian} (Wald)
distribution,
\begin{equation}
  \Tfin_i \;\sim\; \IG(\mu_i,\lambda_i),
  \qquad
  \mu_i \;=\; \frac{L}{v_i},
  \qquad
  \lambda_i \;=\; \frac{L^2}{2\Ddiff_i},
  \label{eq:wald-params}
\end{equation}
with density, for $t>0$,
\begin{equation}
  f_{\IG}(t;\mu,\lambda)
   \;=\; \sqrt{\frac{\lambda}{2\pi t^{3}}}\;
   \exp\!\left(-\,\frac{\lambda\,(t-\mu)^2}{2\mu^{2}\,t}\right).
  \label{eq:ig-density}
\end{equation}

Two derivation checks make the parameterisation transparent. First, the mean: the expected
finishing time is the distance divided by the drift,
$\E[\Tfin_i] = \mu_i = L/v_i$, exactly the deterministic ``run the trip at speed $v_i$'' answer.
Second, the variance. For the Inverse-Gaussian, $\Var(\Tfin_i)=\mu_i^3/\lambda_i$, so
substituting \eqref{eq:wald-params},
\begin{equation}
  \Var(\Tfin_i)
   \;=\; \frac{\mu_i^3}{\lambda_i}
   \;=\; \frac{(L/v_i)^3}{L^2/(2\Ddiff_i)}
   \;=\; \frac{2\Ddiff_i\,L}{v_i^{3}}.
  \label{eq:wald-var}
\end{equation}
The variance of the finishing time is set by the diffusion $\Ddiff_i$ and \emph{grows linearly in
the trip length} $L$. This single fact has large downstream consequences (Section~\ref{sec:cvlaw}
and Chapter~\ref{ch:theory}): it predicts, with \emph{no free parameter}, that longer races are
intrinsically more predictable.

\begin{remark}[Why first-passage, not additive Gaussian]
One could model $\Tfin_i$ directly as $\mu_i + \varepsilon_i$ with Gaussian $\varepsilon_i$. That
is rejected on physical grounds: finishing times are strictly positive and right-skewed (a horse
can lose far more time to trouble than it can ever gain), and the first-passage construction
delivers exactly that --- the Inverse-Gaussian is supported on $(0,\infty)$ and right-skewed,
with the skew arising mechanically from the absorbing boundary. Empirically the additive-Gaussian
``\code{time}'' surrogate scores worse than the Wald law, and the Wald law ties or beats the
log-normal (velocity-multiplicative) surrogate while resting on a single physical diffusion
constant.
\end{remark}

\section{The three knobs}
\label{sec:threeknobs}

Equations \eqref{eq:wald-params}--\eqref{eq:wald-var} expose precisely three places a covariate
can act, and the entire feature programme is organised around them.

\begin{enumerate}[label=\textbf{Knob \arabic*.},leftmargin=*]
  \item \textbf{The drift $v_i$ (aerobic ceiling $\vc$).} Through $\mu_i=L/v_i$, anything that
        raises the sustainable speed lowers the expected time. This is the channel of raw ability:
        form, rating, sectional critical velocity, opponent-adjusted collateral form, speed
        figures, barrier trials, mass. The drift is built up in
        Chapters~\ref{ch:pop}--\ref{ch:competition} and assembled in Chapter~\ref{ch:assembly}.
  \item \textbf{The reserve $\Wp$ (a late-race boost to the effective drift).} The anaerobic
        reserve shortens the effective trip the aerobic engine must cover: a horse with reserve
        $\Wp$ behaves as if running a distance $L$ but with an effective drift scaled up by
        $(1+\Wp/L)$, so that
        \begin{equation}
          \mu_i \;=\; \frac{\Leff}{v_i\,\bigl(1+\Wp_i/L\bigr)},
          \label{eq:mu-wprime}
        \end{equation}
        where $\Leff$ is the geometric length corrected for ground-loss (Chapter~\ref{ch:geometry}).
        The factor $(1+\Wp/L)$ is the production form of Hill's critical-velocity law
        $T=(L-\Wp)/v_c$ (Chapter~\ref{ch:energetics}): the reserve buys back distance, and
        because it is divided by $L$ it \emph{bites harder over sprints than over staying trips},
        which is exactly the distance-aware behaviour the physiology demands. Reserve phenomena
        --- late-race fade resistance, the positional-fade cost a front-runner pays for setting
        the pace --- load this knob (Chapter~\ref{ch:assembly}, the ``T1'' recoupling).
  \item \textbf{The diffusion $\Ddiff_i$ (consistency / racing luck).} Through
        $\lambda_i=L^2/2\Ddiff_i$, the diffusion sets the spread of a horse's time around its
        mean. Because the winner is an \emph{order statistic}, variance is as load-bearing as the
        mean: a high-$\Ddiff$ horse has a fat left tail and upsets more often than its drift
        implies; a metronomic horse upsets less. This is the one knob the model exploited last and
        least; the draw$\to\Ddiff$ trip-variance term (Chapter~\ref{ch:geometry}) and the
        surface-dependent diffusion are its current uses.
\end{enumerate}

\begin{principle}[Right signal, right knob]
A feature is correctly structured only when it loads the knob its mechanism acts on. Loading a
reserve phenomenon (late fade) onto the aerobic ceiling $\vc$ is a \emph{correlated
approximation} --- it raises or lowers the whole velocity curve instead of just the closing kick
--- and is strictly dominated, where the data can resolve it, by loading $\Wp$. Likewise a
consistency signal belongs on $\Ddiff$, not on the mean.
\end{principle}

\section{The closed-form Plackett--Luce surrogate}
\label{sec:pl}

For very wide parameter sweeps it is convenient to have a sampling-free win-probability map. If
the finishing-time noise is taken to be Gumbel rather than Inverse-Gaussian, the argmin in
\eqref{eq:argmin} has the closed-form multinomial-logit (Plackett--Luce) solution. Writing the
deterministic mean times $\mu_i$ and a scale $\sigma_t$, the win probability is the softmax of the
negative mean times,
\begin{equation}
  p_i \;=\; \frac{\exp(-\mu_i/\sigma_t)}{\sum_{j}\exp(-\mu_j/\sigma_t)} ,
  \qquad \mu_i = \frac{\Leff}{v_i(1+\Wp_i/L)} .
  \label{eq:pl}
\end{equation}
Plackett--Luce additionally hands back the \emph{entire} finishing-order distribution in closed
form (the probability of any permutation factorises as a sequence of softmaxes over the remaining
field), which makes it the natural probe for place and exotic structure. It is, however, the
\emph{free / non-interacting} theory: independent finishing times, no traffic, no draft, no
tactical coupling. The production model uses the Wald Monte-Carlo (Chapter~\ref{ch:assembly}); PL
is retained as the instant scout and as the conceptual baseline against which the interaction
programme of Chapter~\ref{ch:theory} is defined.

\begin{remark}[The surrogate is faithful for win, by construction]
Interactions (draw, mass, traffic) have repeatedly tested as low-leverage for the \emph{win}
order statistic specifically --- the knife-edge for first place is dominated by the front of the
field, where the mean-field treatment is adequate. The pure-\code{numpy} Wald/PL surrogates
therefore match the full interacting Monte-Carlo on win to within noise while running roughly
$35\times$ faster, which is what makes the diagnose/recalibrate loop of Chapter~\ref{ch:method}
tractable. The interacting simulator is reserved for the place/exotic regime, where whole-field
ordering --- and hence the many-body terms --- matters.
\end{remark}

\section{The distance law of variance}
\label{sec:cvlaw}

Combine the mean $\mu_i=L/v_i$ and the variance \eqref{eq:wald-var} into the coefficient of
variation of the finishing time:
\begin{equation}
  \mathrm{CV}(\Tfin_i)
   \;=\; \frac{\sqrt{\Var(\Tfin_i)}}{\E[\Tfin_i]}
   \;=\; \frac{\sqrt{2\Ddiff_i L/v_i^3}}{L/v_i}
   \;=\; \sqrt{\frac{2\Ddiff_i}{L\,v_i}}
   \;\propto\; \frac{1}{\sqrt{L}} .
  \label{eq:cvlaw}
\end{equation}
The relative spread of finishing times \emph{falls as one over the square root of the trip
length}. Favourites should win more, and upsets should be rarer, over longer races --- a
prediction that contains no fitted parameter, falling out purely from the geometry of
first-passage. Chapter~\ref{ch:theory} develops the refinement (a distance-independent variance
floor injected by the start and by traffic, which would make sprints \emph{excess}-dispersed
relative to \eqref{eq:cvlaw}) and the order-statistic corollary that the $\Ddiff$ knob has more
leverage in larger fields. This prediction is, at the time of writing, under empirical validation
(queue item \#33).

\section{Summary}

The production win model is the map
\[
  \underbrace{(\vc_i,\ \Wp_i,\ \Ddiff_i)_{i=1}^N}_{\text{per-horse mechanism}}
   \;\longmapsto\;
  \Tfin_i\sim\IG\!\Big(\tfrac{\Leff_i}{\vc_i(1+\Wp_i/L)},\ \tfrac{L^2}{2\Ddiff_i}\Big)
   \;\longmapsto\;
  p_i = \Pr[\Tfin_i = \min_j \Tfin_j].
\]
The remainder of the document is devoted to the first arrow: how each of the three per-horse
parameters is built, leak-free, from data and mechanism. The second arrow is the Wald law of this
chapter; the third is the Monte-Carlo argmin of Chapter~\ref{ch:assembly}.

\chapter{Recovering the latents: critical velocity and anaerobic reserve}
\label{ch:energetics}

\section{The critical-power model of a running horse}

The drift $v_i$ and the reserve $\Wp_i$ of Chapter~\ref{ch:gen} are not free numbers; they are the
two parameters of a \emph{critical-power} physiology, the standard mechanistic model of endurance
performance imported from human exercise science. A horse is idealised as possessing
\begin{itemize}
  \item a \textbf{critical velocity} $\vc$ (m/s): the fastest pace it can hold purely
        \emph{aerobically}, i.e.\ indefinitely on the timescale of a race; and
  \item a finite \textbf{anaerobic reserve} $\Wp$: a fixed quantity of work it can spend running
        \emph{above} $\vc$, after which it must drop back to (or below) $\vc$.
\end{itemize}
The defining relation of the critical-power model is hyperbolic: the distance $L$ a horse can
cover in time $T$ satisfies $L = \vc\,T + \Wp$, i.e.\ it sustains $\vc$ for the whole trip and the
reserve $\Wp$ buys an extra block of distance. Solving for time gives \emph{Hill's
critical-velocity law},
\begin{equation}
  T \;=\; \frac{L - \Wp}{\vc},
  \label{eq:hill}
\end{equation}
which is exactly the deterministic skeleton behind the production mean
\eqref{eq:mu-wprime}: dividing numerator and denominator by $L$,
$T = \frac{L}{\vc}\bigl(1 - \Wp/L\bigr) \approx \frac{L}{\vc(1+\Wp/L)^{-1}}$, and the production
code carries the reserve as the multiplicative late boost $(1+\Wp/L)$ on the drift, which is the
same first-order effect with a numerically gentler, strictly-positive form. The units used
throughout are metres, seconds, and m/s; $\Wp$ is carried as the work proxy
$\int (v-\vc)^{+}\,dt$ with units $\mathrm{m^2/s}$, which is monotone in the true energetic
reserve and suffices for a relative model.

\section{The inverse problem: $(\vc,\Wp)$ from one run's sectionals}

HKJC publishes \emph{sectional} split times: a race is divided into segments, and the time to
cover each segment is recorded. By HKJC's verified convention the final segments are each
$\SI{400}{m}$ and the first segment carries the variable remainder, so for a race of length $L$
split into $n$ sections the per-section distances are
\begin{equation}
  d_1 = L - 400(n-1), \qquad d_2=\dots=d_n = 400,
  \label{eq:sectiondist}
\end{equation}
and the per-section velocities are $v_k = d_k / t_k$ from the published split times $t_k$. The
velocity profile $(v_1,\dots,v_n)$ is the raw kinematic signature of a single run, and from it we
recover that run's $(\vc,\Wp)$ by inverting the critical-power picture.

\paragraph{Estimating $\vc$.} The critical velocity is the speed the horse \emph{settles to}
once the reserve is spent --- the late-race sustainable pace. It is estimated robustly as the
slower of the final two section velocities,
\begin{equation}
  \widehat{\vc} \;=\; \min(v_{n-1},\, v_{n}),
  \label{eq:vc-est}
\end{equation}
the $\min$ guarding against a final-section sprint (a horse with reserve still in hand kicks on
the line, which would over-read $\vc$). Typical recovered values sit around
$\widehat{\vc}\approx \SI{16.8}{m/s}$.

\paragraph{Estimating $\Wp$.} The reserve is the work run \emph{above} that critical velocity,
accumulated over the sections where the horse exceeded it:
\begin{equation}
  \widehat{\Wp} \;=\; \sum_{k=1}^{n} \bigl(v_k - \widehat{\vc}\bigr)^{+}\, t_k,
  \label{eq:wp-est}
\end{equation}
where $(\cdot)^{+}=\max(\cdot,0)$. Two auxiliary descriptors of the same profile are retained for
the trajectory-shape features of Chapter~\ref{ch:features}: the peak section velocity
$v_{\text{peak}}=\max_k v_k$ (the anaerobic ceiling) and the fade
$v_{\text{peak}}/v_n - 1$ (how far the horse dropped from its peak by the line).

\section{Why per-run estimates are noisy --- and must be pooled}

The estimates \eqref{eq:vc-est}--\eqref{eq:wp-est} are \emph{per-run} and individually noisy, for
reasons that are physical rather than statistical artefacts:
\begin{itemize}
  \item an easy winner is eased down through the line, depressing $v_{n}$ and hence under-reading
        $\widehat{\vc}$;
  \item a horse caught in traffic runs a distorted profile that reflects the race shape, not its
        capacity;
  \item the pace of the race --- which the horse does not choose alone --- sets how much of the
        reserve is visible.
\end{itemize}
The latent $(\vc,\Wp)$ are properties of the \emph{horse}, approximately stable across its
career; any single run is a noisy glimpse of them through the lens of that day's tactics. This is
precisely the structure a hierarchical model is built for, and it motivates
Chapter~\ref{ch:pop}: rather than trust one run, we pool a horse's runs toward a population that
is itself conditioned on the horse's class and trip, with the degree of pooling set by how much
evidence the horse has provided. The per-run estimator of this chapter is the \emph{measurement};
the population model of the next chapter is the \emph{estimator of the latent} from those
measurements.

\begin{remark}[A note on what $\vc$ ``is'' in the production model]
The production drift is built around $\vc$ as the organising latent, but $\vc$ should be read as
a \emph{calibrated relative ability in velocity units}, not as a literal physiological constant
measured to three significant figures. Its value is anchored by the sectional estimator above, but
it is then shrunk, rating-conditioned, and nudged by the ability channels of
Chapters~\ref{ch:features}--\ref{ch:competition}, all of which live in the same m/s units so that
they add coherently. The physiology fixes the \emph{structure} (a sustainable ceiling plus a
finite reserve, distance-aware through $\Wp/L$); the calibration fixes the \emph{numbers}.
\end{remark}

\section{A generative race model from $(\vc,\Wp)$ alone}

Before the full stochastic apparatus of Chapter~\ref{ch:gen}, it is worth recording the simplest
generative model the energetics already support, because it shows the mechanism cleanly. Given
per-horse $(\vc_i,\Wp_i)$ and a trip $L$, Hill's law \eqref{eq:hill} gives a deterministic
predicted time $\widehat T_i = (L-\Wp_i)/\vc_i$; adding race-day noise $\varepsilon_i$ and ranking
yields a finishing order,
\begin{equation}
  \Tfin_i \;=\; \frac{L-\Wp_i}{\vc_i} + \varepsilon_i,
  \qquad
  \text{winner} = \argmin_i \Tfin_i .
  \label{eq:naive-gen}
\end{equation}
This is dynamics, not regression: change the distance, or a horse's $(\vc_i,\Wp_i)$, and the
whole finishing-order distribution moves mechanistically --- a horse with a large reserve gains
relatively more over a sprint than over a staying trip, automatically, because $\Wp$ enters as
$\Wp/L$. The production model replaces the additive-Gaussian noise of \eqref{eq:naive-gen} with
the physically-correct first-passage (Wald) noise of Chapter~\ref{ch:gen} --- strictly positive,
right-skewed, with a single diffusion constant and the $\Var\propto L$ law --- but the skeleton is
this critical-power generator.

\chapter{The ability population: hierarchical empirical Bayes}
\label{ch:pop}

\section{The shrinkage problem}

Chapter~\ref{ch:energetics} produced, for each run of each horse, a noisy pair
$(\widehat\vc,\widehat\Wp)$. We need a per-horse \emph{estimate} of the latent ability to feed the
drift, and it must satisfy three requirements simultaneously: it must (i) borrow strength across
horses so a lightly-raced horse is not at the mercy of one fluke run, (ii) condition on the
horse's class and trip so the prior it shrinks toward is appropriate, and (iii) be strictly
leak-free for walk-forward backtesting. The classical tool is a one-way Normal--Normal
random-effects model fitted by empirical Bayes, with a rating-conditioned fixed effect; this
chapter develops it for each of $\vc$ and $\Wp$ independently.

\section{The Normal--Normal random-effects model}

Work with one parameter at a time (say $\vc$; $\Wp$ is identical). Let $y_{hr}$ denote the per-run
estimate \eqref{eq:vc-est} for the $r$-th run of horse $h$. Model it as a fixed effect plus a
horse random effect plus within-horse noise:
\begin{equation}
  y_{hr} \;=\; \underbrace{f(b_{hr}, R_{hr})}_{\text{fixed effect}}
            \;+\; \underbrace{\theta_h}_{\text{horse effect}}
            \;+\; \underbrace{\varepsilon_{hr}}_{\text{within-horse noise}},
  \qquad
  \theta_h \sim \Norm(0,\tau^2), \quad \varepsilon_{hr}\sim\Norm(0,\sigma^2).
  \label{eq:randeff}
\end{equation}
Here $b_{hr}\in\{\text{sprint},\text{mile},\text{middle},\text{staying}\}$ is the distance band
of the run and $R_{hr}$ is the horse's official rating going into it. The two variance components
have direct interpretations: $\sigma^2$ is the \emph{within-horse} run-to-run variance (how noisy
a single sectional estimate is), and $\tau^2$ is the \emph{between-horse} variance (how much
genuine spread there is in true ability after conditioning on class and trip).

\subsection{The rating-conditioned fixed effect}

The fixed effect is not a single grand mean but a per-band intercept plus a global rating slope,
\begin{equation}
  f(b,R) \;=\; \beta_b \;+\; \gamma\,(R-\bar R),
  \label{eq:fixed}
\end{equation}
so that a horse's prior centres on the value expected \emph{for its class and trip}, not on the
global average. The band intercepts $\beta_b$ capture that $\vc$ falls with distance (a sprinter's
sustainable velocity exceeds a stayer's); the slope $\gamma$ captures that a higher-rated horse
has, on average, a higher critical velocity. The slope is estimated by an ordinary least-squares
fit of the band-centred residual on the centred rating,
\begin{equation}
  \widehat\gamma
   \;=\; \frac{\sum_{hr}(R_{hr}-\bar R)\,(y_{hr}-\beta_{b_{hr}})}{\sum_{hr}(R_{hr}-\bar R)^2},
  \label{eq:slope}
\end{equation}
after which the rating-explained part is removed to form the pure horse-plus-noise residual
$u_{hr} = y_{hr} - \beta_{b_{hr}} - \widehat\gamma(R_{hr}-\bar R)$. Conditioning on rating is what
makes the population a \emph{segmented-by-class} prior rather than a mono-population; it matters
because a debut or lightly-raced horse shrinks almost entirely to this prior, so the prior had
better know the horse's class.

\subsection{Estimating the variance components (ANOVA method of moments)}

The two variance components are estimated by the ANOVA method of moments --- robust, closed-form,
and assumption-light. The within-horse variance is the pooled sample variance over horses with at
least two runs,
\begin{equation}
  \widehat{\sigma}^2
   \;=\; \frac{\sum_{h:\,n_h\ge 2}(n_h-1)\,s_h^2}{\sum_{h:\,n_h\ge 2}(n_h-1)},
  \qquad
  s_h^2 = \frac{1}{n_h-1}\sum_{r}(u_{hr}-\bar u_{h})^2,
  \label{eq:sigma2}
\end{equation}
where $n_h$ is horse $h$'s run count and $\bar u_h$ its mean residual. The between-horse component
follows from equating the between-group mean square to its expectation. With
$\mathrm{SSB}=\sum_h n_h \bar u_h^2$ (the residuals are already fixed-effect-centred so the grand
mean is zero), $H$ horses, and total runs $n_\bullet=\sum_h n_h$,
\begin{equation}
  \mathrm{MSB} = \frac{\mathrm{SSB}}{H-1},
  \qquad
  n_0 = \frac{1}{H-1}\Bigl(n_\bullet - \frac{\sum_h n_h^2}{n_\bullet}\Bigr),
  \qquad
  \widehat{\tau}^2 = \max\!\Bigl(\frac{\mathrm{MSB}-\widehat\sigma^2}{n_0},\ \epsilon\Bigr),
  \label{eq:tau2}
\end{equation}
where $n_0$ is the usual effective group size for unbalanced designs and the floor $\epsilon$
keeps $\widehat\tau^2$ strictly positive. (The implementation, \code{calibration.py}, computes
exactly these quantities in vectorised \code{polars}/\code{numpy}.)

\section{The partially-pooled per-horse estimate}

With the components in hand, the empirical-Bayes posterior mean for horse $h$ is the
fixed-effect prediction plus a \emph{shrunken} version of its own mean residual,
\begin{equation}
  \boxed{\;\widehat{\vc}_h \;=\; f(b,R) \;+\; \lambda(n_h)\,\bar u_h,\qquad
  \lambda(n_h) \;=\; \frac{\tau^2}{\tau^2 + \sigma^2/n_h}\;}
  \label{eq:shrink}
\end{equation}
The shrinkage weight $\lambda(n_h)\in[0,1)$ is the heart of the model. Read its limits:
\begin{itemize}
  \item $n_h\to 0$ (debut): $\lambda\to 0$, the estimate is the pure rating-conditioned population
        mean --- the horse has provided no evidence, so it is its class.
  \item $n_h\to\infty$ (established): $\lambda\to 1$, the estimate trusts the horse's own runs
        almost entirely.
  \item Intermediate $n_h$: $\lambda$ interpolates, and the cross-over is governed by the
        \emph{noise-to-signal ratio} $\sigma^2/\tau^2$: a noisier measurement (large $\sigma^2$)
        or a more homogeneous population (small $\tau^2$) means more pooling for the same run
        count.
\end{itemize}
The posterior also has a variance, used later by the precision-fusion of Chapter~\ref{ch:fusion}:
\begin{equation}
  \mathrm{sd}(\widehat{\vc}_h)
   \;=\; \Bigl(\frac{n_h}{\sigma^2} + \frac{1}{\tau^2}\Bigr)^{-1/2},
  \label{eq:postsd}
\end{equation}
the standard ``precisions add'' result: the posterior precision is the prior precision $1/\tau^2$
plus $n_h$ run-precisions $1/\sigma^2$. This identity --- that each independent source of evidence
contributes its precision additively --- is the conceptual seed of the entire fusion architecture.

\section{Leak-free variants: leave-one-out and strictly-prior}

The shrinkage formula \eqref{eq:shrink} must be made leak-free in two distinct settings, and the
implementation provides both.

\paragraph{Leave-one-out (for in-sample calibration targets).} When the energetics of a run are
used to \emph{calibrate} against that same run's outcome, the run must be excluded from its own
horse's pool. The LOO estimate replaces $\bar u_h$ and $n_h$ by their values with the current run
dropped,
\begin{equation}
  \bar u_{h}^{(-r)} \;=\; \frac{\sum_{r'\ne r} u_{hr'}}{n_h-1},
  \qquad
  \lambda(n_h-1),
  \label{eq:loo}
\end{equation}
computed in vectorised form as (group sum $-$ this run)/(group count $-1$). This is the quantity a
race model may consume without a run's own outcome leaking into the fit that predicts it.

\paragraph{Strictly-prior expanding window (for walk-forward forecasting).} For a genuine
forecast of a race on date $d$, the estimate may use only runs \emph{strictly before} $d$. The
trailing variant sorts each horse's runs by date and forms an expanding mean over earlier runs
only,
\begin{equation}
  \widehat{\vc}_h(d) \;=\; f(b,R) \;+\; \lambda(n_{<d})\,
   \frac{1}{n_{<d}}\!\!\sum_{r:\,\text{date}(r)<d}\!\! u_{hr},
  \label{eq:trailing}
\end{equation}
with $n_{<d}$ the count of prior runs (and a fallback to the pure population mean when
$n_{<d}=0$). The population parameters $(\beta_b,\gamma,\sigma^2,\tau^2)$ are themselves fit on the
training window only, so the whole pipeline --- population and per-horse estimate --- is leak-free.
This trailing estimate is the one the walk-forward backtest of Chapter~\ref{ch:method} actually
consumes.

\section{Output of the population layer}

The population layer delivers, for every runner in every race, a leak-free
$(\widehat\vc_h,\widehat\Wp_h)$ together with the horse's prior-run count $n_h$ (which the
shrinkage and the later fusion both need) and the population's variance components. These are the
\emph{baseline} drift and reserve --- the ``rating-conditioned mean plus shrunken own-history''
that a horse carries into a race before any of the feature channels of the next chapters nudge it.
In the assembled drift of Chapter~\ref{ch:assembly} this appears as the leading terms
$\code{pop.v\_c.fixed}(b,R) + \texttt{SHRINK\_MULT}\cdot\lambda(n)\,u_{\vc}$ and the analogous
reserve term, on top of which every ability channel adds.

\chapter{The feature layer: leak-free ability channels}
\label{ch:features}

\section{Role of the feature layer}

The population layer (Chapter~\ref{ch:pop}) gives each horse a baseline drift and reserve from its
own sectional history shrunk toward a rating-conditioned prior. The feature layer supplies the
\emph{additional} signals that move a horse's drift away from that baseline: measurements of
ability, condition, intent, and style that the sectional history does not already capture. Every
feature in this chapter obeys the leak-free invariant of Chapter~\ref{ch:intro} --- it is a
function of a horse's strictly-prior runs (or of pre-race public boundary conditions) only --- and
every feature is interpreted as loading one of the three knobs. This chapter catalogues the
channels; Chapter~\ref{ch:fusion} explains the principled way to \emph{combine} the velocity-unit
ability channels, and Chapter~\ref{ch:assembly} writes the assembled drift in full.

A unifying remark frames what follows. The drift's ability channels --- rating, form, sectional
$\vc$, collateral form, barrier trial, speed figure, allometric mass --- are all \emph{noisy
measurements of the same latent}, the critical velocity $\vc$. The principled combination is a
precision-weighted posterior mean,
\begin{equation}
  \vc \;=\; \frac{\sum_k \pi_k\,\vc^{(k)}}{\sum_k \pi_k},
  \qquad \pi_k = 1/\sigma_k^2 ,
  \label{eq:precision-fusion}
\end{equation}
and the production drift $\vc^{\mathrm{eff}} = \vc^{\text{base}} + \sum_k g_k\, s_k$ is the
first-order linearisation of \eqref{eq:precision-fusion}, in which each gain $g_k$ is the relative
precision of channel $k$. This is why the gains are \emph{calibrated, not free}: a channel that
is collinear with one the market already prices adds no precision and earns a gain near zero, and
over-coupling a channel past its true precision \emph{hurts}. Keep \eqref{eq:precision-fusion} in
view; it is the logic behind every gain in the constants table.

\section{Form: the dominant ability signal}

\subsection{Career-mean relative finish}

The single strongest feature is \emph{form} --- a horse's history of finishing well relative to
its fields. The atom is the within-race \emph{relative finish}, which places a result on a
$[0,1]$ scale invariant to field size,
\begin{equation}
  \mathrm{rf} \;=\; 1 - \frac{\text{finish\_pos}-1}{\text{field\_size}-1}
  \quad\in[0,1], \qquad \mathrm{rf}=1 \text{ for a win.}
  \label{eq:relfin}
\end{equation}
The career-mean form is the strictly-prior expanding mean of $\mathrm{rf}$ over a horse's earlier
runs (all venues count),
\begin{equation}
  \mathrm{form}_h(t) \;=\;
   \frac{\sum_{s\le t}\mathrm{rf}_s - \mathrm{rf}_t}{n_{\le t}-1},
  \label{eq:form}
\end{equation}
the by-one-shifted expanding mean that excludes the current run; it is null on debut.

\subsection{Recency weighting: the biggest single lever found}

The career mean is stale: over the long careers typical in Hong Kong, a horse's form a dozen runs
ago says little about today. Replacing the flat expanding mean by an \emph{exponentially-weighted}
one --- recent runs weighted more --- is the largest single improvement in the model's history
(roughly $+7.8$ percentage points of market). With a half-life $h$ measured in \emph{races}
(not days), the recency form is
\begin{equation}
  \mathrm{form}^{\mathrm{rec}}_h(t)
   \;=\; \frac{\sum_{s<t} w_{t-s}\,\mathrm{rf}_s}{\sum_{s<t} w_{t-s}},
   \qquad w_{\Delta} = 2^{-\Delta/h},
  \label{eq:recform}
\end{equation}
again strictly prior (the exponentially-weighted mean is shifted by one). The half-life is baked
at $h=\texttt{FORM\_HALF\_LIFE}=2.0$ races: a clean optimum, with the career mean ($h\to\infty$)
the worst case. Two mechanisms explain why recency wins: it tracks the horse's \emph{current}
ability through the form cycle, and it matches how the market itself weights recent evidence.
Form enters the drift with a per-venue gain $\texttt{FORM\_GAIN\_BY\_VENUE}$ (ST $5.0$, HV $4.0$;
HV form is more weakly coupled because it was found to be over-weighted there), centred at the
neutral value $0.5$:
\begin{equation}
  \Delta\vc^{\text{form}}_i \;=\; \texttt{FORM\_GAIN}\cdot\bigl(\mathrm{form}^{\mathrm{rec}}_i - 0.5\bigr).
  \label{eq:formcoupling}
\end{equation}

\subsection{Prior-start count}

A simple but structurally important companion is the strictly-prior count of official starts,
$\text{prior\_starts}_h(t) = n_{\le t}-1$. A value of zero marks a \emph{debut} --- the rows for
which the barrier-trial channel (Section~\ref{sec:trial}) is the only on-track form available, and
the rows where the fusion of Chapter~\ref{ch:fusion} does its most distinctive work.

\section{Jockey and trainer channels}

\subsection{Jockey strike rate}

The booked jockey carries a trip-quality signal (skill at pace judgement, path-finding,
rail-saving) that is partly a real $\vc$ effect and partly an information proxy (fancied horses
attract the best riders). It is measured as the jockey's strictly-prior expanding win strike rate,
\begin{equation}
  \mathrm{jky\_sr}_j(t) \;=\;
   \frac{\sum_{s\le t}\mathbb{1}[\text{win}]_s - \mathbb{1}[\text{win}]_t}{n^{(j)}_{\le t}-1},
  \label{eq:jky}
\end{equation}
keyed to each horse by the jockey riding it, and coupled with gain
$\texttt{JOCKEY\_GAIN}=6.0$ centred on the field mean. Because the jockey signal is partly
endogenous to the very market we fit, it is kept as a measurement but not treated as a clean single
mechanism.

\subsection{Trainer strike rate and the hot-streak overlay}

The trainer analogue $\mathrm{trn\_sr}$ is the stable's strictly-prior expanding win rate, a
career-level competence signal that turns out to be largely redundant with jockey and rating (it
adds $\approx 0.2$pp and is parked). The sharper, \emph{recency}-weighted version --- the stable's
\emph{hot streak} --- does convert, because it carries \emph{intent}: a yard suddenly winning
races is placing horses to win, a pre-race public signal the market prices. It is an
exponentially-weighted mean of the trainer's win indicator over its strictly-prior runners, with a
half-life of $30$ runs,
\begin{equation}
  \mathrm{trn\_hot}_T(t) \;=\;
   \frac{\sum_{s<t} 2^{-(t-s)/30}\,\mathbb{1}[\text{win}]_s}{\sum_{s<t} 2^{-(t-s)/30}},
  \label{eq:trnhot}
\end{equation}
within-race centred and coupled with $\texttt{TRN\_HOT\_GAIN\_BY\_VENUE}$ (ST $2.0$, HV $1.0$).

\section{Running-style and positional traits}

Two distinct notions of ``running style'' are computed, and the distinction matters.

\paragraph{Speed-based style.} $\mathrm{style}$ is the strictly-prior expanding mean of the
normalised early \emph{relative velocity} $\mathrm{rel\_early}$ (front-runner high, closer low),
from the sectional data. It is retained as a feature but is \emph{not} what the production
positional term rides on.

\paragraph{Position-based early position.} $\mathrm{early\_pos}$ is the strictly-prior expanding
mean of the normalised early field \emph{position} $\mathrm{early\_pos\_norm}$ (1 = led early, 0 =
settled last) from the run-policy dataset. Winners average $\approx 0.95$ against the
$\approx 0.76$ field mean, so high values mean habitually nearer the front early. The two style
notions correlate only $\approx 0.12$ --- position and speed are genuinely different traits ---
and it is the positional one that feeds the production model, through the \emph{positional-fade}
mechanism: the non-interacting drift--diffusion (Chapter~\ref{ch:gen}) over-rates habitual
front-runners, because it never charges them for the anaerobic reserve they burn setting a
contested pace with no cover, and symmetrically under-rates closers. The correction is to debit a
front-runner's reserve, which in the production model loads the $\Wp$ knob
(Chapter~\ref{ch:assembly}); $\mathrm{early\_pos}$ is the leak-free trait that drives it. This is
the measured, mean-field shadow of the pace equilibrium developed in Chapter~\ref{ch:theory}.

\section{Velocity speed figures}

Rank-based form \eqref{eq:relfin} censors the winning margin --- a neck and ten lengths both
score $\mathrm{rf}=1$. A complementary, kinematically-grounded channel is the \emph{speed figure},
a horse's average velocity expressed relative to a condition par so that runs across different
trips and going are comparable. With $v_{\text{run}} = L/\Tfin$ (m/s) and a par velocity
$\mathrm{par}(\text{going},L)$ taken as the mean run velocity over all races at that
going-and-distance,
\begin{equation}
  \mathrm{speed\_fig} \;=\; v_{\text{run}} - \mathrm{par}(\text{going},L),
  \label{eq:speedfig}
\end{equation}
then the strictly-prior expanding mean per horse. Because $\mathrm{speed\_fig}$ is already in m/s,
it couples to $\vc$ with a kinematic coefficient $\approx 1$ rather than a free gain --- a clean
example of a velocity-unit channel. The par is a race-condition constant (not the horse's own
outcome), so computing it over the full panel is leak-free. In practice the rank-form channel
carries most of the win-relevant ability information (its pace/class invariance is an advantage),
and the speed figure is largely redundant for win, but it is the template for the trial and
collateral channels that are \emph{not}.

\section{Barrier-trial debut form}
\label{sec:trial}

A barrier trial is an official, timed practice gallop --- crucially, it is a horse's \emph{only}
on-track form before its racing debut, and a fitness read on a return from a layoff. For debutants
the market is otherwise efficient with information the model never sees; the trial closes that gap
and was the first new-data lever to convert (ST $+1.9$pp). The trial speed figure mirrors
\eqref{eq:speedfig} on trial data: $v_{\text{trial}} = L/\text{trial\_time}$, a par
$\mathrm{par}(\text{going},L)$ over all trials at that condition, the difference
$\mathrm{trial\_fig} = v_{\text{trial}}-\mathrm{par}$, and the strictly-prior expanding mean per
horse. Trials happen on mornings, races on race-days, so a horse's trial and race timelines never
share a date and a backward as-of join is leak-free. A subtlety of the join key is recorded in the
code: trial data carry the long registration id, but the model panels key on the short stable
brand, so the table emits the brand aliased to \code{horse\_id}. The additive coupling is
$\texttt{TRIAL\_GAIN\_BY\_VENUE}$ (ST $1.5$, HV $1.0$); Chapter~\ref{ch:fusion} replaces this
additive form by the precision-fusion, which is its principled successor.

\section{Trajectory-shape traits}

The sectional $\vc$/speed-figure channels read a run's \emph{average} energetics; the full
trajectory dataset carries the per-section velocity \emph{profile}, exposing the \emph{shape} of
how a horse runs --- mechanics the average hides. Two physically-grounded traits are extracted per
run and made strictly-prior:
\begin{align}
  \mathrm{sustain} &= \frac{\bar v_{\text{sections}}}{v_{\text{peak}}}\ \in(0,1],
    && \text{sustained-speed ability (1 = flat cruise)};\\
  \mathrm{fade} &= v_{\text{last}} - v_{\text{peak}}\ \le 0,
    && \text{fade resistance (0 = no late deceleration)}.
  \label{eq:trajtraits}
\end{align}
The fade is the kinematic signature of the reserve $\Wp$ emptying in the closing stages, so it is
a \emph{reserve} phenomenon and belongs on the $\Wp$ knob. Each raw trait is within-race
standardised (a contemporaneous condition normaliser, not an outcome) and then averaged
strictly-prior per horse. Sustained-speed tests as already-priced (parked); fade resistance
converts as a small reserve term ($\texttt{FADE\_W\_GAIN}=8.0$ on $\Wp$, Chapter~\ref{ch:assembly}).

\section{Rating-trajectory process features}

The official handicap rating is a \emph{time series} per horse, not just a latest point, and its
\emph{trajectory} encodes the career cycle the point misses. From the dated rating series the model
derives, per update, the velocity $\mathrm{rtg\_vel}=\Delta\mathrm{rtg}/k$, the acceleration
$\mathrm{rtg\_accel}$, the maturity count $\mathrm{rtg\_nupd}$, and --- the one that converts ---
the \emph{peak gap}
\begin{equation}
  \mathrm{rtg\_peakgap} \;=\; \mathrm{rtg}(t) - \max_{s\le t}\mathrm{rtg}(s) \ \le 0,
  \label{eq:peakgap}
\end{equation}
how far the horse sits below its own career-best rating. This is a \emph{process} feature: it
carries an \emph{absolute career landmark} the point rating cannot express. A horse declining off
its peak runs below its (lagging) official rating, and the market discounts this; the model
corrects it with a negative gain $\texttt{PEAKGAP\_GAIN\_BY\_VENUE}$ (ST $-0.015$, HV $-0.025$),
within-race centred. The rating velocity and acceleration do \emph{not} convert --- they merely
restate recent form, which the form channel already holds --- a clean illustration that a process
feature pays only where it adds information orthogonal to the point and to form. The trajectory is
attached by a strict backward as-of join with \code{allow\_exact\_matches=False}, so a race sees
only rating updates strictly before it.

\section{Summary of channels and their knobs}

Table~\ref{tab:channels} collects the channels, their measured signal, and the knob each loads.
The recurring discipline is visible in the rightmost column: \emph{ability} signals load $\vc$,
\emph{reserve} signals load $\Wp$, and \emph{consistency} signals load $\Ddiff$ --- and the next
chapter shows how the $\vc$ channels should be fused rather than merely summed.

\begin{table}[h]
\centering
\small
\begin{tabular}{@{}lll@{}}
\toprule
Channel & Measured signal & Knob \\
\midrule
Recency form & exp-weighted prior relative finish & $\vc$ (drift) \\
Sectional $\vc$ residual & shrunk own-history critical velocity & $\vc$ (drift) \\
Jockey / trainer-hot & prior \& recency win strike rate & $\vc$ (drift), part info \\
Speed figure & prior mean $v_{\text{run}}-\mathrm{par}$ (m/s) & $\vc$ (drift) \\
Barrier trial & prior mean $v_{\text{trial}}-\mathrm{par}$ (m/s) & $\vc$ (drift, debut) \\
Collateral form & opponent-adjusted Massey velocity (m/s) & $\vc$ (drift) \\
Allometric mass & $(m/\bar m)^{b-1}$ & $\vc$ (drift, ST) \\
Rating peak-gap & own-peak deficit ($\le 0$) & $\vc$ (drift, process) \\
Fade resistance & late-vs-peak section velocity & $\Wp$ (reserve) \\
Positional fade & habitual early position & $\Wp$ (reserve) \\
Draw $\times$ field crowding & arc-length ground loss & $\Leff$ (drift, geometry) \\
Draw $\times$ bend $\times$ extremity & lateral exposure & $\Ddiff$ (variance, ST) \\
Surface (AWT) & dirt kickback & $\Ddiff$ (variance) \\
\bottomrule
\end{tabular}
\caption{The feature channels, their leak-free signal, and the generative knob each loads.}
\label{tab:channels}
\end{table}

\chapter{Coherent precision-fusion of velocity observations}
\label{ch:fusion}

\section{Why additive coupling is only the linearisation}

The drift's ability channels are noisy measurements of one latent, $\vc$, and the exact Bayesian
combination is the precision-weighted posterior \eqref{eq:precision-fusion}. The additive
production drift is its first-order linearisation, adequate when every channel has a stable
precision and a horse has a typical amount of evidence. It is \emph{not} adequate when a channel's
relative informativeness \emph{changes with the amount of other evidence the horse has} --- which
is exactly the situation for a barrier trial. A trial is enormously informative about a
\emph{debutant} (it is the only on-track data that exists) and nearly worthless about a horse with
forty races behind it (those races dominate). A single additive gain cannot express ``weight this
a lot at debut and let it fade as races accumulate''. The coherent fusion does precisely that, by
returning to the precision algebra of Chapter~\ref{ch:pop} and adding the trial as its own
evidence channel with its own precision.

\section{The fusion of an observation channel into the $\vc$ posterior}

Recall the posterior precision identity \eqref{eq:postsd}: precisions add. Decompose the evidence
about a horse's $\vc$ residual into three independent sources, each contributing its own
precision:
\begin{equation}
  \pi_0 = \frac{1}{\tau^2}\ \text{(population prior)},\qquad
  \pi_{\text{race}} = \frac{1}{\sigma^2}\ \text{(one race obs)},\qquad
  \pi_{\text{trial}} = \rho\,\pi_{\text{race}}\ \text{(one trial obs)} ,
  \label{eq:precisions}
\end{equation}
where $\rho = \pi_{\text{trial}}/\pi_{\text{race}}$ is the \emph{relative} precision of a trial
observation versus a race observation --- the one new number the fusion introduces, baked as
$\texttt{TRIAL\_PI\_RATIO}$. A horse carrying $n$ prior races and $m$ prior trials has a posterior
$\vc$-residual that weights the trial evidence by
\begin{equation}
  \boxed{\;
  w_{\text{trial}}(n,m)
   \;=\; \frac{m\,\pi_{\text{trial}}}{\pi_0 + n\,\pi_{\text{race}} + m\,\pi_{\text{trial}}}
   \;=\; \frac{m\rho\,\pi_{\text{race}}}{\pi_0 + n\,\pi_{\text{race}} + m\rho\,\pi_{\text{race}}}\; }
  \label{eq:wtrial}
\end{equation}
This is the fraction of total posterior precision that the trial channel supplies. Read its
behaviour: as the race count $n$ grows, the denominator grows and $w_{\text{trial}}\to 0$ --- the
trial fades exactly as the horse's racing record accumulates, with no threshold or hand-tuned
schedule. At debut ($n=0$, $m=1$) it takes its maximal value
$w_0 = \rho\pi_{\text{race}}/(\pi_0+\rho\pi_{\text{race}})$.

\subsection{Normalising to the additive coupling at debut}

The fusion is designed so that for a debutant it reproduces the (calibrated, converting) additive
trial coupling of Section~\ref{sec:trial}, and only \emph{deviates} from it by fading as the horse
races. This is achieved by normalising $w_{\text{trial}}$ to unity at debut and multiplying the
additive trial term by the normalised weight,
\begin{equation}
  \widehat w(n,m) \;=\; \frac{w_{\text{trial}}(n,m)}{w_0},
  \qquad
  \Delta\vc^{\text{trial}}_i \;=\; \widehat w(n_i,m_i)\cdot \texttt{TRIAL\_GAIN}\cdot \mathrm{trial\_fig}_i .
  \label{eq:trialterm}
\end{equation}
At debut $\widehat w = 1$ and the fusion equals the additive coupling; as the horse races
$\widehat w$ decays smoothly toward $0$. When the fusion is active ($\rho>0$) the plain additive
trial term is switched off, so the two are mutually exclusive structural choices for the same
signal. The result: the trial enters at \emph{its own} noise level, recovering the full additive
signal for a debutant (ST $+6$pp on debut races specifically) and the calibration benefit of a
denser prior for lightly-raced horses, while not contaminating the estimate of an established
horse.

\subsection{The structural lesson behind the fusion}

The fusion's design encodes a hard-won negative result. An earlier attempt folded the barrier
trial into the \emph{sectional-$\vc$ residual} channel, which the production model globally shrinks
by $\texttt{SHRINK\_MULT}\approx 0.5$ (a multiplier calibrated to the \emph{noisy} sectional
estimator). That diluted a comparatively \emph{high}-precision signal by forcing it through a
shrinkage meant for a low-precision one, and it underperformed the direct additive coupling for the
win order statistic --- even though it did tighten the ability estimate (improve ECE) for
lightly-raced horses, confirming the denser-prior mechanism was real. The lesson is structural,
not a refutation of the mechanism:

\begin{principle}[Each observation type keeps its own precision]
A fused observation must enter at its \emph{own} estimated precision. A shared shrinkage --- making
a trial inherit the sectional channel's variance --- is the wrong structure. The correct fusion
assigns each channel (races, trials, gallops, race \emph{phases}) its own precision component, so
each enters at its own noise level. ``A feature failing under one coupling means \emph{that
structure} is wrong, not that the mechanism is dead.''
\end{principle}

\section{Recovering the trial critical velocity on the race scale}

For the fusion to be coherent, the trial's signal must be on the \emph{same physical scale} as the
race $\vc$ residual, which requires recovering a per-horse critical velocity from a trial rather
than just a relative figure. Trials are run in small batches that share a leader; the recovery
proceeds in two steps:
\begin{enumerate}[label=(\roman*)]
  \item the batch leader's section split profile is passed through the \emph{same}
        \code{estimate\_run\_energetics} machinery as race energetics (Chapter~\ref{ch:energetics})
        to get a batch critical velocity $\vc^{\text{batch}}$ --- the settled late velocity the
        field fades to;
  \item each horse ran the same trial at its own \code{trial\_time}, so its critical velocity
        scales with its average pace relative to the batch leader,
        \begin{equation}
          \vc^{\text{trial}}_i \;=\; \vc^{\text{batch}}\cdot
           \frac{\text{batch\_time}}{\text{trial\_time}_i}.
          \label{eq:trialvc}
        \end{equation}
\end{enumerate}
By construction the resulting $\vc^{\text{trial}}$ has a mean near the race $\vc$ mean
($\approx \SI{16.82}{m/s}$), so a single train-only affine map suffices to place it on the
race-\emph{residual} scale (the map is fit on training rows only, hence leak-free). The table then
emits the strictly-prior expanding mean $\mathrm{trial\_vc}$ and the prior trial count
$\mathrm{trial\_m}=m$, the two quantities \eqref{eq:wtrial} consumes.

\section{Extending the fusion to training gallops}
\label{sec:gallop}

The architecture was built explicitly as the \emph{substrate for a third channel}: routine
training gallops. A gallop is the densest but noisiest velocity observation of a horse --- there
are tens of thousands of them (the \code{trackwork.parquet} panel holds $30\text{k}+$ Banister
gallop sessions), they happen continuously between races, but a solo morning gallop pace is not a
contested race pace and carries far less information per observation than a trial. The coherent
fusion accommodates this exactly by giving the gallop channel its own (low) precision ratio,
\begin{equation}
  \pi_{\text{gallop}} = \rho_g\,\pi_{\text{race}},\qquad
  w_{\text{gallop}}(n,m,m_g)
   = \frac{m_g\,\pi_{\text{gallop}}}
          {\pi_0 + n\,\pi_{\text{race}} + m\,\pi_{\text{trial}} + m_g\,\pi_{\text{gallop}}},
  \label{eq:gallopfusion}
\end{equation}
with $\rho_g$ small (gallops are noisy) and baked as $\texttt{GALLOP\_PI\_RATIO}$. This is an
active work item (queue \#30): it requires a \code{gallop\_vc\_table} mirroring
\eqref{eq:trialvc} --- careful, because a solo gallop has no opponents to set a par --- and a full
$\sim\!1.5\text{k}$-horse backfill of the join. The hope, grounded in the precision algebra, is
that the denser gallop prior converts where the sparse trial axis alone does not, particularly at
Happy Valley. The same template generalises once more in Chapter~\ref{ch:theory}: a race's
\emph{closing} phase is a higher-precision $\vc$ channel than its trip-average, so ``give each
\emph{phase} its own precision'' is the temporal analogue of ``give each \emph{observation type}
its own precision''.

\begin{remark}[Fusion is the right home for all velocity-unit channels]
The collateral-form rating (Chapter~\ref{ch:competition}), the speed figure, and the finishing
sectional are all m/s measurements of $\vc$ and, in the limit of the programme, belong in one
coherent shrinkage posterior over all velocity observations, each down-weighted by its own noise,
with the correct nonlinear $\lambda(n_{\text{eff}})$. The current production drift realises this as
a calibrated linear combination plus the explicit trial fusion; the unifying target is the single
posterior of \eqref{eq:precision-fusion} made exact.
\end{remark}

\chapter{The competition axis: relational collateral form}
\label{ch:competition}

\section{Why ability is relational}

The hardest and most under-exploited dimension of the problem is \emph{competition}. The handicap
system actively compresses the field --- the handicapper assigns weight to engineer a blanket
finish --- which is precisely why win probabilities are close and why beating the market is hard.
The consequence for modelling is that \emph{raw} ability is the wrong object: beating a strong
field should count for more than the same performance against a weak one. The win model, at
equilibrium, \emph{is itself a rating system}; the right ability measure is the consistent
solution of an opponent-adjusted system.

The first attempt at this axis --- a scalar Elo / relative-form rating --- failed, and the failure
is instructive. A scalar rating is collinear ($\approx 0.70$) with the official rating the market
already prices, and it discards the only thing competition actually does, which is to make outcomes
\emph{relational}: a horse's effective ability is a \emph{function of the opponent set}, not a
context-free number. A front-runner among closers steals an uncontested lead and arrives with
reserve unspent; the same front-runner among other front-runners is forced into a contested duel
that burns $\Wp$ against the air with no draft, and fades. Identical horse, different field,
different result. A scalar cannot represent a quantity that depends on who else is in the race. The
correct object is \emph{relational}, and the relational treatment converts strongly (HV $+2.3$pp,
with ECE improving) where the scalar one vanished --- the canonical illustration that ``a dead
first attempt is not a dead axis.''

\section{The Massey collateral-form rating}

\subsection{Construction}

The collateral-form rating is an opponent-adjusted \emph{velocity} rating in m/s, built as an
online Massey least-squares system so that it is strictly leak-free. The atom is again the
condition-normalised speed figure of a run, $s_i = L/\Tfin_i - \mathrm{par}(\text{going},L)$ (m/s,
the same kinematic figure as \eqref{eq:speedfig}). Races are processed in date order. The
\emph{field strength} of a race $R$ is the mean \emph{prior} collateral rating $\bar R$ of its
runners (unrated newcomers count as neutral $0$). A runner's observation from race $R$ combines its
\emph{margin against the field it actually ran} with the \emph{strength of that field}:
\begin{equation}
  \mathrm{obs}_i \;=\; \underbrace{\bigl(s_i - \overline{s}_R\bigr)}_{\text{margin vs this field}}
   \;+\; \kappa\,\underbrace{\bar R}_{\text{strength of this field}},
  \label{eq:collat-obs}
\end{equation}
and the horse's rating is the expanding mean of its $\mathrm{obs}$ over its strictly-prior runs.
The $\kappa\bar R$ term is the \emph{collateral propagation}: a margin posted against a strong
field inherits that field's strength, which is one Gauss--Seidel sweep of the Massey
least-squares system in date order. The least-squares value is $\kappa=1$; $\kappa<1$ damps the
propagation. Leak-freeness is preserved by emitting the rating a horse carries \emph{into} a race
(the running per-horse mean over strictly-prior runs) and only \emph{then} updating that running
mean with the current race's observation --- the same strictly-prior guarantee as the form table,
and the par is a condition constant rather than the horse's own outcome, so the full panel is safe.

\subsection{Coupling}

The emitted $\mathrm{collat\_form}$ is within-race centred (a field-mean shift cancels in the win
probability) and coupled to the drift with a per-venue gain $\texttt{COLLAT\_GAIN\_BY\_VENUE}$ (ST
$1.0$, HV $2.5$). The much larger HV gain is consistent with HV's tighter, more compressed fields
making opponent adjustment more informative there; the collateral channel closes roughly a third
of the ST--HV skill gap.

\section{Deepening the rating (v2)}

The production ST collateral rating uses a refined construction (\code{collateral\_form\_table\_v2})
that pushes the Massey solution closer to its true fixed point and conditions it more carefully.
All refinements preserve the strictly-prior / online leak-free discipline. The baked ST
construction is $\{\,n_{\text{sweeps}}=3,\ \text{half\_life}=4.0,\ \text{band\_par}=\text{True}\,\}$;
HV keeps the single-sweep table (band conditioning fragments its tight trip range). The four
refinement axes:

\begin{enumerate}[label=\textbf{R\arabic*.},leftmargin=*]
  \item \textbf{Deeper propagation ($n_{\text{sweeps}}>1$).} A single online pass propagates
        opponent strength only one hop along the collateral chain. Additional \emph{Jacobi}
        fixed-point sweeps re-solve the strictly-prior race graph before each new race-date is
        emitted, each sweep recomputing every prior race's observation using the previous sweep's
        \emph{frozen} ratings as the field-strength source. Strength then propagates one more hop
        per sweep, toward the true Massey fixed point. Crucially the extra sweeps touch only
        strictly-prior races, so the rating emitted for a race still contains no information from
        it or any later race. Baked at three sweeps on ST.
  \item \textbf{Recency weighting (\emph{half\_life}).} The per-horse rating can be an
        exponentially-weighted mean of its observations rather than a flat expanding mean, so a
        stale opponent rating decays --- the recency-form lever (Section \ref{ch:features}) applied
        to the competition channel. With decay $\delta = 2^{-1/\text{half\_life}}$ the running
        numerator and denominator are updated as $\delta\cdot(\text{old}) + (\text{new})$. Baked at
        a four-race half-life on ST.
  \item \textbf{Nonlinear margin (\emph{margin\_pow}) and strength-of-schedule asymmetry
        (\emph{asym}).} The within-race margin can be raised to a sign-preserving power
        ($\mathrm{sign}(m)\,|m|^{p}$, $p>1$ rewards big margins, $p<1$ compresses), and the field
        strength inherited on a \emph{loss} can be weighted differently from a \emph{win} (kept
        term $\kappa\bar R$ for a beat, $\kappa\cdot\mathrm{asym}\cdot\bar R$ for a loss), encoding
        that losing to a strong field tells you less than beating one. Both default to the linear,
        symmetric baked values.
  \item \textbf{Trip-band conditioning (\emph{band\_par}).} The par velocity and the margin are
        computed within (going, distance-\emph{band}), and the field-strength inheritance is
        conditioned on the trip band, so a 1200\,m margin feeds a horse's rating only for nearby
        trips. Ability is partly trip-specific; band conditioning enforces like-for-like
        comparison. A horse carries a per-band rating with a fallback to its global mean. Baked on
        ST.
\end{enumerate}

The deepened ST construction added $+1.3$pp \emph{on top of} the trial fusion (super-additive),
and was what broke ST through $65\%$; its ECE also improves. This is the competition axis,
operationalised --- the relational rating the ``win model is a rating system at equilibrium''
framing predicts, built jointly with the leak-free discipline rather than bolted on.

\begin{remark}[The free theory and the interacting theory]
The production Wald/PL surrogate is the \emph{free}, non-interacting theory: context-free drifts,
no traffic, no draft. The collateral rating injects the \emph{first relational moment} of
competition --- opponent strength --- into the per-horse drift, which is why it survives the
collinearity that killed the scalar. The deeper relational structure --- how a horse's effective
ability depends on the \emph{style composition} of \emph{this} field --- is the interacting theory
of Chapter~\ref{ch:theory}, not yet in production.
\end{remark}

\chapter{The arena: geometry, draw, and ground loss}
\label{ch:geometry}

\section{Track geometry and the movable rail}

The draw bias is fundamentally \emph{geometric}: on a turn, an inner lane is a shorter path, so a
low-drawn horse that holds the rail runs less ground. The model carries an explicit geometry,
parameterised by \emph{venue $\times$ surface $\times$ course}, because HKJC turf tracks use a
\emph{movable running rail} --- the race's \code{Course} (A, A+2, B, B+2, B+3, C, C+3) shifts the
rail outward to spread wear, changing the usable width and slightly the home-straight length. Three
structural regimes are encoded as \code{TrackGeometry} objects:
\begin{itemize}
  \item \textbf{Happy Valley} (turf, $\SI{1450}{m}$ circumference, right-handed): a tight,
        bend-dominated circuit ($\approx 470 + 530\,\si{m}$ of turn within a $\SI{1450}{m}$ lap)
        --- the strongest draw bias in Hong Kong.
  \item \textbf{Sha Tin turf} (galloping oval, $\SI{1900}{m}$, right-handed, $\SI{430}{m}$ home
        straight): long straights and two $\approx \SI{495}{m}$ bends, hence a \emph{weaker}
        geometric draw bias than HV.
  \item \textbf{Sha Tin straight $\SI{1000}{m}$ chute}: a separate course with \emph{no bends}, so
        no bend-geometry draw bias at all (only going/watering gradients, deliberately unmodelled).
  \item \textbf{Sha Tin all-weather} (dirt, $\SI{1560}{m}$, fixed single rail, uniform
        $\SI{22.8}{m}$ width).
\end{itemize}
The geometry exposes a \emph{bend fraction} --- the fraction of a race of length $L$ run on turns,
computed by sampling positions along the trip and testing each against the bend segments,
\begin{equation}
  \mathrm{bend\_fraction}(L) \;=\; \frac{1}{P}\sum_{p=1}^{P}\mathbb{1}\bigl[\text{position } d_p
   \text{ is on a turn}\bigr],
  \label{eq:bendfrac}
\end{equation}
which is the draw-bias exposure of the race. It is $\approx 0$ on the straight chute (which is why
the draw term is zeroed there) and largest at HV.

\section{Draw as ground loss on the drift}

A horse drawn $\mathrm{Dr}$ positions off the rail runs a longer arc round the turns. The
production model represents this as an \emph{effective lengthening} of the trip the drift must
cover,
\begin{equation}
  \Leff \;=\; L\,\bigl(1 + g_{\text{draw}}\,(\mathrm{Dr}-1)\bigr),
  \label{eq:leff}
\end{equation}
which feeds straight into the Wald mean \eqref{eq:mu-wprime}: a wider draw means a longer effective
trip, a longer first-passage time, and a lower win probability. The gain $g_{\text{draw}}$ is
\emph{per-venue and geometric}, baked as $\texttt{DRAW\_GAIN\_BY\_VENUE}$ (ST $0.003$, HV $0.008$):
HV, bend-dominated, wants roughly $2.7\times$ the ST penalty. On the all-weather surface the rail
bias is stronger and is overridden by $\texttt{DRAW\_GAIN\_AWT}=0.008$. On any course with
$\mathrm{bend\_fraction}<0.05$ (the straight chute) the gain is set to zero --- no arc, no penalty.

\begin{remark}[A lesson in always sweeping]
An early diagnosis read the HV draw residual as ``over-weighted'' and proposed setting the gain to
zero; the sweep proved the \emph{opposite} --- zeroing it craters HV to $41\%$, because HV was in
fact massively \emph{under}-penalised by the old global $0.003$. The episode is the canonical
reminder that the sign of a residual coupling must be established by a sweep against the production
path, never assumed.
\end{remark}

\section{Draw $\times$ field size: the crowding interaction}

The geometric penalty is amplified when the field is crowded: a wide draw costs more in a full
field (more bodies to be trapped behind, fewer gaps to cross to the rail) than in a sparse one ---
the first \emph{interaction} moment of the draw mechanism. The production form scales the draw gain
by a field-size factor centred at twelve runners, clamped so the penalty never reverses in a small
field:
\begin{equation}
  g_{\text{draw}}^{\text{eff}} \;=\; g_{\text{draw}}\cdot
   \max\!\bigl(0,\ 1 + g_{\text{fs}}\,(N-12)\bigr),
  \label{eq:drawfs}
\end{equation}
with $\texttt{DRAW\_FS\_GAIN\_BY\_VENUE}$ (HV $0.25$, ST $0.0$). It converts only at HV (the
tightest track, where collision physics dominates: HV $+0.3$pp, flat at ST). It is explicitly an
\emph{interim stopgap} --- a scalar shadow of a genuinely two-dimensional geometry$\times$control
interaction (Chapter~\ref{ch:theory}) --- and is flagged as such pending the joint treatment.

\section{Draw as variance: loading the diffusion knob}

The deeper insight is that draw acts not only on the \emph{mean} (the arc-length drift penalty) but
on the \emph{variance}: a wide-drawn horse in a big field has a more uncertain trip --- it may get
across and save ground, or it may be hung wide and buried --- and that uncertainty is a property of
$\Ddiff$, not of the drift. This is the first deliberate use of the third knob. Each horse's
diffusion is inflated by a draw-derived ``trip risk'',
\begin{equation}
  \Ddiff_i \;=\; D_{\text{eff}}\,\bigl(1 + g_{\text{dvar}}\cdot
   \mathrm{bend\_fraction}(L)\cdot \mathrm{extremity}_i\bigr),
  \qquad
  \mathrm{extremity}_i = \frac{2\,\bigl|\mathrm{Dr}_i - \tfrac{N+1}{2}\bigr|}{N-1},
  \label{eq:drawvar}
\end{equation}
where $\mathrm{extremity}_i\in[0,1]$ is the field-normalised lateral extremity of the draw (a
rail or widest-out horse is extreme, a midfield draw is not), and the risk is gated by the bend
fraction (no bends, no trip risk). A larger $\Ddiff_i$ means a larger Wald variance
\eqref{eq:wald-var} and so more upsets from that horse. The gain
$\texttt{DRAW\_VAR\_GAIN\_BY\_VENUE}$ is ST $0.5$, HV $0.0$.

The ST-only split is \emph{predicted by theory}, a clean theory-confirms-data check. The winner is
the left-tail extreme of $N$ finishing-time draws; in a larger field the win requires a deeper tail
event, so per-horse variance has more leverage on \emph{who wins the larger $N$ is}. ST runs the
big fields, HV the small ones, so the variance knob should pay on ST and barely at all on HV ---
exactly the observed split (see Chapter~\ref{ch:theory}, the extreme-value corollary).

\section{Surface as diffusion}

Surface enters the variance knob directly: the all-weather (dirt) surface is more diffusive than
turf --- kickback and a less uniform surface make times less repeatable --- so its diffusion is
scaled by $\texttt{awt\_diff\_mult}=2.0$,
\begin{equation}
  D_{\text{eff}} \;=\; \texttt{diff\_coef}\cdot
   \begin{cases} 2.0 & \text{AWT}\\ 1.0 & \text{turf,}\end{cases}
   \qquad \texttt{diff\_coef}=170,
  \label{eq:awtdiff}
\end{equation}
which corrects a pronounced AWT over-confidence (raising the model's AWT win-probability dispersion
from an over-tight $12.7\%$ toward a calibrated $23.7\%$). The empirical distance-law study of
Chapter~\ref{ch:state} independently recovers an AWT diffusion roughly twice the turf value,
vindicating this multiplier from data.

\section{What geometry does and does not yet capture}

The terms above --- $\Leff$ on the drift, the crowding amplifier, the draw$\to\Ddiff$ variance, and
the surface diffusion --- are the \emph{mean-field shadows} of draw: per-horse scalars that capture
the field-averaged, first-moment consequences of geometry. They do \emph{not} capture the genuine
two-dimensional, many-body phenomenon in which the cost of a draw is a functional of the jockey's
control policy and of the style composition of the neighbouring horses. That is the
geometry$\times$control game of Chapter~\ref{ch:theory}, the single most promising piece of open
theory, for which the baked draw terms are explicit interim stopgaps and for which the
computer-vision trajectory programme is being built to supply the missing lateral coordinate.

\chapter{Assembly: the production drift, reserve, and predict loop}
\label{ch:assembly}

\section{The assembled effective critical velocity}

Collecting every $\vc$-loading channel from Chapters~\ref{ch:pop}--\ref{ch:competition}, the
effective critical velocity each runner $i$ carries into the simulator is
\begin{align}
  \vc^{\mathrm{eff}}_i \;=\;\;
   & \underbrace{f_{\vc}(b, R_i)}_{\text{rating-conditioned pop.\ mean}}
   \;+\; \underbrace{\texttt{SHRINK\_MULT}\cdot\lambda_{\vc}(n_i)\,u_{\vc,i}}_{\text{shrunk own-history residual}}
   \;+\; \underbrace{\widehat w(n_i,m_i)\,g_{\text{trial}}\,\mathrm{trial\_fig}_i}_{\text{trial fusion (Ch.\,\ref{ch:fusion})}}
   \nonumber\\[1mm]
   &\;+\; g_{\text{form}}\,\bigl(\mathrm{form}_i - 0.5\bigr)
   \;+\; g_{\text{jky}}\,\bigl(\mathrm{jky}_i - \overline{\mathrm{jky}}\bigr)
   \;+\; g_{\text{collat}}\,\widetilde{\mathrm{cfm}}_i
   \nonumber\\[1mm]
   &\;+\; g_{\text{pgap}}\,\widetilde{\mathrm{pgap}}_i
   \;+\; g_{\text{thot}}\,\widetilde{\mathrm{thot}}_i
   \;+\; (\text{speed-fig, trainer, sustain, additive-trial terms, mostly off}),
  \label{eq:vc-eff}
\end{align}
where a tilde denotes within-race centring (a constant field shift cancels in the win
probability), $b$ is the distance band, $R_i$ the rating, $n_i$ the prior-race count, $m_i$ the
prior-trial count, and the gains are the per-venue baked constants of Appendix~\ref{app:constants}.
The mass channel multiplies the whole drift allometrically (Section~\ref{sec:allo}). Equation
\eqref{eq:vc-eff} is the linearised precision-weighted posterior \eqref{eq:precision-fusion}:
every term is a channel's signal times its relative precision, and a channel collinear with one
the market already prices contributes a gain near zero.

\section{The assembled reserve}

The anaerobic reserve carries the late-race / intertemporal phenomena. It is the
rating-conditioned reserve mean plus the shrunk own-history reserve residual, plus the two
reserve couplings moved onto this knob (the ``T1'' recoupling):
\begin{equation}
  \Wp^{\mathrm{eff}}_i \;=\;
   f_{\Wp}(b,R_i)
   \;+\; \texttt{W\_PRIME\_GAIN}\cdot\lambda_{\Wp}(n_i)\,u_{\Wp,i}
   \;+\; \underbrace{\texttt{FADE\_W\_GAIN}\cdot\mathrm{fade}_i}_{>0:\ \text{holds speed}\Rightarrow\text{more reserve}}
   \;+\; \underbrace{\texttt{POSFADE\_W\_GAIN}\cdot\widetilde{\mathrm{epos}}_i}_{<0:\ \text{front-runner burns}\,\Wp}.
  \label{eq:wp-eff}
\end{equation}
The signs are mechanistic: a horse that holds its speed late ($\mathrm{fade}>0$ in the trait's
sign convention) has demonstrated reserve, so $\texttt{FADE\_W\_GAIN}=+8.0$; a habitual
front-runner (high $\mathrm{early\_pos}$) burns reserve setting the pace, so
$\texttt{POSFADE\_W\_GAIN}=-80.0$. Placing these on $\Wp$ rather than $\vc$ is the
right-signal-right-knob principle in action: their effect becomes \emph{distance-aware for free}
through the $\Wp/L$ form of the late boost \eqref{eq:mu-wprime}, biting harder over sprints --- and
indeed the gain from the recoupling lands in HV sprints ($+0.3$pp), the $\Wp/L$ signature. The
predecessors $\texttt{FADE\_GAIN}$ and $\texttt{POS\_FADE\_GAIN}$ (which loaded $\vc$) are retired
to zero.

\section{Allometric mass}
\label{sec:allo}

Body mass enters through the power-balance derivation of critical velocity. With aerobic power
scaling allometrically as $P_c = \kappa m^{b}$ and the steady-state balance
$\vc = P_c/(K_{\text{LOCO}} m g)$, the sustainable velocity scales as
\begin{equation}
  \vc \;\propto\; m^{\,b-1},
  \label{eq:allo}
\end{equation}
so the drift is multiplied by the field-relative factor $(m_i/\bar m)^{b-1}$, where $\bar m$ is the
geometric-mean mass of the field and the declared weight is pre-race (hence leak-free). The
exponent is one \emph{venue-invariant} physiological constant; it is applied where it converts
(ST, $b=1.1$) and is off where other channels already price it (HV, $b=1.0$). That mass is
low-leverage at HV is a statement about \emph{collinearity with already-priced channels}, not a
claim that HV horses obey a different physiology --- a distinction the first-principle discipline
insists on.

\section{The predict loop}

For each race in the test set the model:
\begin{enumerate}[label=\arabic*.]
  \item resolves the per-venue gains, the geometry \code{track\_for(venue, distance, surface)}, and
        the per-race precision ratio $\rho$;
  \item assembles $\vc^{\mathrm{eff}}_i$ and $\Wp^{\mathrm{eff}}_i$ from \eqref{eq:vc-eff} and
        \eqref{eq:wp-eff};
  \item forms the Wald parameters with the geometry of Chapter~\ref{ch:geometry},
        \begin{equation}
          \mu_i = \frac{\Leff_i}{\vc^{\mathrm{eff}}_i\,(1+\Wp^{\mathrm{eff}}_i/L)},
          \qquad
          \lambda_i = \frac{L^2}{2\,\Ddiff_i};
        \end{equation}
  \item draws $S$ Monte-Carlo finishing times per runner from the Inverse-Gaussian,
        $t^{(s)}_i \sim \IG(\mu_i,\lambda_i)$, with a \emph{race-deterministic seed} so the noise
        is identical across model variants (the basis of the paired conversion test of
        Chapter~\ref{ch:method});
  \item counts argmin wins across the $S$ samples and Laplace-smooths to a probability,
        \begin{equation}
          p_i \;=\; \frac{\#\{s: t^{(s)}_i = \min_j t^{(s)}_j\} + 1}{S + N},
          \label{eq:laplace}
        \end{equation}
        the $+1/+N$ guaranteeing every runner a strictly positive probability (no $\log 0$ in the
        scoring).
\end{enumerate}
The default sample count is $S=\texttt{N\_SAMPLES}=6000$. The Plackett--Luce branch
\eqref{eq:pl} substitutes a closed-form softmax for steps 4--5 when an instant, sampling-free probe
is wanted.

\section{The constants are the calibration}

Every gain in \eqref{eq:vc-eff}--\eqref{eq:wp-eff} and every geometry constant lives in the single
module \code{hkjc.model.constants}, each a mechanism-coupling gain set by the conversion test of
Chapter~\ref{ch:method}, not a free regressor weight. Many are per-venue dictionaries, because the
two tracks are structurally different arenas (HV tight and bend-dominated with small fields; ST a
galloping oval with big fields), and a few are per-surface. Appendix~\ref{app:constants} reproduces
the full table with values and mechanisms. The discipline that keeps this from being an
overfit pile of knobs is the subject of the next chapter: each constant must be both
physics-derivable \emph{and} move the out-of-sample needle, confirmed by a paired fixed-seed test
on both venues.

\chapter{The theoretical frontier: the race as a game}
\label{ch:theory}

The production model is the \emph{free}, non-interacting theory: independent first-passage times
with mean-field corrections. The remaining skill gap to the market lives in what that theory omits
--- the \emph{interactions} between horses, which are products across dimensions, not main effects,
and which a purely additive model cannot represent. This chapter sets out the theory of those
interactions along three axes --- longitudinal (pace), lateral (geometry), and temporal (phase) ---
each formulated as a differential / mean-field game. It is the source of the forward research
queue.

\section{The longitudinal pace game}
\label{sec:pacegame}

Fix a trip $L$. Each horse $i$ carries a reserve $\Wp_i$ above a critical velocity $\vc_i$, and its
jockey chooses an effort/pace policy $\pi_i(t)$ --- a velocity profile, equivalently a
reserve-depletion schedule --- to minimise its own finishing time
$T_i = \int_0^L ds/v_i(t)$ subject to the reserve budget $\int (v_i-\vc_i)^{+}dt \le \Wp_i/k$ and
the position dynamics $\dot x_i = v_i$. The payoff is a \emph{tournament}: only finishing
\emph{order} pays. Because each horse's optimal effort is a best response to the others' (the value
of leading depends on who contests the lead; the draft you get depends on who sits where), the
field's tactics form a \emph{Nash equilibrium}. With many runners the rivals enter only through the
aggregate \emph{pace}, so the large-field limit is a \emph{mean-field game}: each horse best-responds
to the field's pace distribution, which is itself the aggregate of those responses --- a fixed point
in pace-policy space.

\paragraph{Two levels.} Above the tactical pace game sits the handicapper as a \emph{Stackelberg
leader}: anticipating the equilibrium, the handicapper assigns weight to compress ability and
engineer a blanket finish. Competition thus enters twice --- the handicapper compresses the
\emph{mean} (already priced, hence tight win probabilities and an efficient rating channel), but the
pace equilibrium reshuffles the \emph{conditional} outcome by style composition, which is
\emph{not} fully priced. That residual is the opening.

\paragraph{The positional fade is the equilibrium's mean-field shadow.} The non-interacting drift
over-rates habitual front-runners (never charged for the contested-pace reserve burn) and
under-rates closers (never credited the reserve they conserve). The baked positional-fade term
(now on the $\Wp$ knob, Chapter~\ref{ch:assembly}) --- front horses winning $\approx 1$pp \emph{less}
than the physics predicts, closers $\approx 1$pp \emph{more}, residual $t\approx -7$ at \emph{both}
venues --- is exactly this bias, measured. The shared sign and magnitude across venues is the
signature of a \emph{mean-field constant}: the field-averaged shadow of the pace equilibrium
projected onto one per-horse scalar. The model has captured only the \emph{zeroth moment} (the
field-mean style offset); the equilibrium's \emph{interaction} structure --- how the fade depends on
\emph{who else} is in the race --- is the next layer.

\paragraph{Testable predictions.} All inputs are pre-race and public, hence leak-free. Ranked:
(1) the front-runner penalty \emph{scales with the count of rival front-runners} (pace
contestedness $n_{\text{FR}}$) and vanishes or reverses at $n_{\text{FR}}=1$; (2) a lone-leader
pace-control bonus (a sign flip, robustly expressed as the interaction with $n_{\text{FR}}$, not a
threshold-fragile standalone cell); (3) closers' bonus scales with $n_{\text{FR}}$; (4) a single
pace-pressure scalar per race summarising the early-position distribution, modulating the fade
gain; (5) pairwise matchup / collateral terms (the interacting Plackett--Luce, highest ceiling,
heaviest build). Predictions 1--3 are one mechanism --- the equilibrium's first moment,
fade-scaling-with-contestedness --- from three angles, a single \texttt{contestedness}-modulated
extension of the positional-fade gain. The first probes did not convert as bare scalars, which is
the signal that the effect is geometry$\times$control coupled (next section), not a simple scalar.

\section{The lateral geometry--control game}
\label{sec:geomgame}

Section~\ref{sec:pacegame} is the along-track direction; draw lives in the \emph{across-track}
direction, which the scalar draw terms of Chapter~\ref{ch:geometry} only shadow. The thesis: draw
is not a static geometric covariate but a \emph{boundary condition on the jockey's control policy}.
Give each horse-jockey a 2-D state $s_i = (x_i, y_i, \Wp_i, v_i)$ --- along-track position, lateral
offset from the rail, reserve, speed --- and two controls: the effort $u_i\ge 0$ above the aerobic
ceiling (depleting $\Wp_i$) and the lateral steering rate $\dot y_i$ (the path: hug the rail vs.\
peel out for clear air). The equations of motion are the 2-D version of the power-balance plant,
with a bend ground-loss term $\propto \kappa(x_i)\,v_i^2\,y_i$ and a lateral \emph{exclusion}
constraint (a horse cannot occupy a rival's lane --- a TASEP-like single-file occupancy in each
lateral band).

\paragraph{The draw is the initial lateral condition; its cost is the equilibrium value-gap.} The
start fixes $y_i(0)=y_0(\mathrm{Dr}_i)$. To save ground a wide horse must \emph{cross} --- spend
$\dot y_i<0$ early --- but crossing is contested: the rail is scarce and the gap must open. So a
wide draw imposes a \emph{control dilemma} --- commit reserve early to cross (paying a $\Wp$ tax up
front, risking being shuffled back) or stay wide (conceding arc length but keeping clear air and
reserve). The realised draw cost is the equilibrium value-gap
$\Delta V_i = V_i(\text{rail}) - V_i(\text{wide})$, a \emph{functional} of the curvature, the
field's lateral congestion, and the style composition of the neighbours --- \emph{not} a constant
$g_{\text{draw}}(\mathrm{Dr}-1)$. Identical draw, opposite cost, decided by who is inside you. This
is why the scalar pace/draw probes did not convert, and why the baked draw terms are interim.

\paragraph{Leak-free observables.} The binding constraint is that \emph{no panel stores the lateral
coordinate} $y_i(t)$; the only across-track datum is the starting offset $\mathrm{Dr}$. What exists
is the longitudinal rank trajectory (rank at several call points), the per-section velocity and
position, and a derived run-policy panel (early/late position, position gain, fade, kick fraction,
effort). The lateral path is therefore latent --- inferred from rank trajectory $\times$ draw
$\times$ geometry. The ranked, tractable, leak-free observables this predicts: \textbf{O1}, a
draw$\times$course (curvature/width) modulation of the draw gain (pure pre-race geometry --- but a
scalar-mean version did \emph{not} convert, rejecting that structure, not the mechanism);
\textbf{O2}, draw $\times$ field-pace-makeup (relational congestion built from the early-position
table); \textbf{O3}, a control-quality / ground-saved trait (rank gained per unit of relative
section velocity --- a jockeyship latent, no CV needed); \textbf{O4}, jockey policy-type clustering;
\textbf{O5}, the full 2-D many-body cell (the simulator's domain). O3 (the control share, available
from existing trajectory data) and the computer-vision lateral paths are the open work.

\section{The temporal phase structure}
\label{sec:phases}

The production Wald law treats the trip as a \emph{single homogeneous} drift--diffusion with one
drift and one diffusion. A real race is three dynamically distinct \emph{phases}, and collapsing
them smears signals that live in different phases and load different knobs. Resolving the phases is a
measurement upgrade with no new mechanism --- the kind of lever (like recency form) the theory says
pays.

\paragraph{Phase 1 --- the start transient.} From rest, acceleration-limited, off the gate; not the
steady-state first-passage. It enters as a random initial condition $T_i = T_{0,i} + T_{\text{main},i}$
with a per-horse \emph{mean} (gate speed / break quality, a stable leak-free trait loading the drift)
and a per-horse \emph{variance} (gate temperament, loading $\Ddiff$). This is the ``start not
explicit'' scorecard gap, given a structure. The break trait is \emph{not} early position: it is the
position change \emph{across the break} --- first-call rank adjusted for draw --- and its
wide-draw interaction is sharp (a slow break from a wide gate forfeits both longitudinal and lateral
ground, missing the gap to cross). The convertible experiment is a leak-free trailing
$\mathrm{gate\_break}$ trait (queue \#31), unmodelled today.

\paragraph{Phase 2 --- the cruise.} Tactical: the pace duel and the lateral navigation. Least about
raw ability, most about control --- which is exactly why the mid-race sectionals are the noisiest
$\vc$ channel.

\paragraph{Phase 3 --- the closing.} The reserve empties, every horse is maximally extended, nobody
holds anything back --- the \emph{least tactically-contaminated} read of raw capacity in the race.
The finishing sectional (last $\approx\SI{400}{m}$, normalised for the pace it was run into and the
going) is therefore the \emph{highest-precision} $\vc$+$\Wp$ channel available, higher than the
trip-average speed figure, which dilutes the clean closing signal with the tactical cruise. The
convertible experiment is a finishing-speed figure entered as its own high-precision channel in the
fusion of Chapter~\ref{ch:fusion} (queue \#32) --- ``give each \emph{phase} its own precision,'' the
temporal analogue of the fusion's central lesson.

\paragraph{Control is intertemporal.} The phase where a control choice is \emph{made} differs from
the phase where its cost is \emph{paid}: reserve spent on an early forward position is reserve
missing in the closing kick. This cross-phase debit is the mechanism behind the baked T1 recoupling
(positional fade $\to \Wp$, distance-aware via $\Wp/L$): an early control decision (lead) debited
from the late reserve. By the same logic fade-resistance, positional-fade, and kick-timing are
facets of one intertemporal $\Wp$-deployment policy, resolved by phase --- their shared home is
$\Wp$, and the phase a trait reads from tells you which sign it carries. A homogeneous single-drift
law cannot carry this coupling; phase-resolving the trip exposes it.

\section{The order-statistic corollaries of variance}
\label{sec:evt}

Two corollaries of the distance-law of variance (Section~\ref{sec:cvlaw}) sharpen the prioritisation
of the variance knob.

\paragraph{The distance-law, validated.} The Wald law predicts, parameter-free,
$\mathrm{CV}(T)\propto 1/\sqrt{L}$ --- longer races intrinsically more predictable. An empirical
study (queue \#33) confirms it on turf: the fitted slope of $\log\mathrm{CV}$ on $\log L$ is
$-0.507$ (95\% CI $[-0.84,-0.17]$), and the backed-out diffusion $D$ shows \emph{no} significant
trend with $L$ (slope $+0.12$, CI $[-0.27,+0.50]$), i.e.\ $D$ is constant in $L$ as the simplest
form of the law predicts --- a clean theory-confirms-data check. The one exception is the
all-weather surface, which is both more diffusive ($D_0\approx 0.94$ vs.\ turf $0.70\,\mathrm{m^2/s}$,
consistent with the baked $\texttt{awt\_diff\_mult}=2$) and carries a large distance-\emph{independent}
floor ($\sigma_0\approx\SI{0.46}{s}$, $\approx 36\%$ of AWT-sprint variance). So a distance-dependent
$D(L)$ floor is \emph{not} warranted on turf but is the right structure for AWT (an additive
$\Var \mathrel{+}= \sigma_0^2$).

\paragraph{Field size amplifies the variance knob.} The winner is the left-tail extreme of $N$
finishing-time draws. In a larger field the win requires a \emph{deeper} tail event, so per-horse
variance $\Ddiff$ has \emph{more leverage on who wins the larger $N$ is}: a high-$\Ddiff$ horse's
fat left tail is worth more in a 14-runner ST handicap than in a 6-runner HV field. The Monte-Carlo
captures this exactly (it simulates argmin over the true field), so no new term is needed --- but it
is a sharp prioritisation rule: \emph{the $\Ddiff$ knob pays in big fields and barely at all in
small ones}. This retro-explains the baked draw$\to\Ddiff$ result converting on ST and being null on
HV (ST runs the big fields), and it directs every future variance lever --- per-horse $\Ddiff_i$
from figure variance (queue \#28), gate temperament, the AWT floor --- to ST and AWT first.

\section{Why this is one programme}

The three games share one structure: the production law is a mean-field projection, and the
remaining signal is in the \emph{interactions} --- relational (who else is in the race), geometric
(where on the track), and sequenced (which phase). Each is a differential game whose mean-field
shadow the model already carries as a scalar (positional fade, draw gain, the homogeneous drift),
and each promises its next increment of skill when the interaction itself --- not its shadow --- is
modelled. The substrate, when these axes are taken up jointly, is an interacting Plackett--Luce with
condensed-matter machinery (exclusion processes for blocking, active matter for drafting, mean-field
games for jockey policy) built on the real geometry, validated against the same conversion test that
gates everything else.

\chapter{Method: diagnosis, the conversion test, and leak-freeness}
\label{ch:method}

\section{The calibration loop}

The model is improved by one disciplined loop:
\begin{equation}
  \text{diagnose} \;\longrightarrow\; \text{recalibrate} \;\longrightarrow\;
  \text{bake} \;\longrightarrow\; \text{re-diagnose}.
  \label{eq:loop}
\end{equation}

\paragraph{Diagnose.} Run the model walk-forward and take the signed residual between the market
and the model win probability, both de-vigged and race-normalised,
\begin{equation}
  \mathrm{resid}_i \;=\; p^{\text{market}}_i - p^{\text{model}}_i .
  \label{eq:resid}
\end{equation}
Regress $\mathrm{resid}$ on each observable mechanism with \emph{race-clustered} standard errors
(runners within a race are not independent --- the probabilities sum to one). The sign and the
$|t|$ say which mechanism the model mis-weights and by how much; the largest is the next thing to
ground in physics. The diagnosis script (\code{market\_residual\_diagnosis.py}) is the instrument.

\paragraph{Recalibrate.} Fix the identified miss with a \emph{single mechanism-grounded constant}
or the right structural change --- never a guessed covariate. If a feature fails to convert,
re-derive the structure (the right latent, precision, interaction, or knob) before declaring the
mechanism dead.

\paragraph{Bake.} The constant goes into \code{hkjc.model.constants} only after it passes the
conversion test below. Its one-line what/units lives with the code; its story --- the deltas, the
ECE, the BAKE/DON'T-BAKE decision --- lives in the journal, not in a comment.

\paragraph{Re-diagnose.} Confirm both $\%$ of market \emph{and} ECE improved out-of-sample (not
just in-sample), then find the next miss.

\section{The conversion test is the gate --- not the t-statistic}

A residual diagnosis only \emph{points} at a candidate; the decision to bake is made by the
conversion test, and the distinction is load-bearing.

\begin{principle}[The gate]
A change is baked only if it improves \emph{$\%$ of market (out-of-sample) and calibration (ECE),
both}, confirmed at \emph{full scale} with a \emph{paired fixed-seed} backtest
($\texttt{PYTHONHASHSEED=0}$, so base and candidate share per-race Monte-Carlo seeds and the delta
is pure mechanism, not noise), on \emph{both venues}.
\end{principle}

Two facts force this discipline:
\begin{itemize}
  \item \textbf{A high $|t|$ residual is not a useful lever.} What matters for $\%$ of market is the
        impact on the \emph{winner's} probability, which a large mean residual on midfield runners
        need not move. The draw mechanism sat for a long time as a $t=-26.5$ ``flat ridge'' worth
        nothing on the kinematic baseline, then added $+1.9$pp once re-added to the \emph{production}
        (Wald) drift it had been dropped from --- the residual was real, but its leverage depended
        entirely on the generative path it loaded.
  \item \textbf{The subsampled quick evaluation over-states a knob $\approx 4\times$.} The crowding
        draw term read $+1.0$ to $+1.6$pp on a 150-race \code{quick\_eval} but only $+0.3$pp HV /
        $0.0$pp ST at full scale --- a fourfold overstatement --- and seed salting alone injects
        $\pm 0.3$ to $0.5$pp of run-to-run noise. So \code{quick\_eval} is used \emph{only} to scout
        direction and sign; nothing is baked on it, and the paired full-scale run is mandatory.
\end{itemize}

\paragraph{Why paired and fixed-seed.} The Monte-Carlo argmin \eqref{eq:laplace} injects sampling
noise of order the effect sizes being chased. Seeding each race deterministically and sharing the
seed between base and candidate makes the noise \emph{common} to both, so it differences out and the
measured $\Delta(\%\text{ of market})$ is the mechanism alone. Without pairing, a $+0.3$pp real
effect is invisible under $\pm 0.5$pp independent noise.

\section{Diagnostics that can mislead}

Two diagnostic traps are recorded so they are not re-sprung. First, a going/class \emph{categorical}
residual test is \emph{uninformative} under race-normalisation, because a race-level constant
vanishes when the residuals are centred within each race; condition calibration must instead be
checked by an ECE \emph{within} each going/class bucket. Second, a residual coupling's sign must be
established by a sweep against the current production path, never read off the raw residual sign ---
the HV draw episode (Chapter~\ref{ch:geometry}) is the standing warning.

\section{Leak-freeness, enforced mechanically}

Every number in the loop is only as trustworthy as the leak-free guarantee beneath it
(Chapter~\ref{ch:intro}). The enforcement is mechanical, not a matter of care:
\begin{itemize}
  \item per-horse features are strictly prior --- expanding means shifted by one, or strict as-of
        joins with \code{allow\_exact\_matches=False};
  \item the energetics population and per-horse residuals are fit on the \emph{training} window of
        each walk-forward fold only;
  \item features are built through an \code{AsOfView} wrapper and \code{assert\_no\_lookahead} is
        run on the fit/predict before any backtest number is trusted;
  \item contemporaneous boundary conditions (draw, geometry, field composition, declared weight)
        are pre-race public and are the \emph{only} non-historical inputs allowed.
\end{itemize}
A leak-free path is never weakened to chase a number; a backtest that has seen the future is worse
than no backtest, because it is confidently wrong.

\section{New data and new structure, not more mean-field}

The recurring strategic finding is that the mean-field ``saturation'' the model hit was a
saturation of \emph{additive terms on the existing race panel} --- every per-horse residual was
collinear with the priced ability channel. The flat fundamentals (draw, carried weight, official
rating, crude form) are exhausted; the market prices the main effects. Gains come only from
(i) measuring a latent better (recency form $\gg$ career mean), (ii) adding a missing mechanism to
the generative law (bend ground-loss in the drift), (iii) bringing in new data the market prices but
the model never saw (barrier trials, training gallops), or (iv) a structural / relational signal the
additive panel cannot express (opponent-adjusted collateral form, the career-landmark process
feature). All are mechanism or measurement upgrades, never re-weightings --- exactly what the
first-principle discipline predicts, and what the 2026 push (Chapter~\ref{ch:state}) proved.

\chapter{Current state, results, and the open queue}
\label{ch:state}

\section{Headline fit}

At the time of writing the production model (\code{SIM=wald}, per-venue calibration) reaches
\begin{center}
  \textbf{ST $\approx 65.0\%$ of market} \qquad \textbf{HV $\approx 59.7\%$ of market},
\end{center}
with a model ECE of $\approx 0.006$, up from $\approx 48\%$ at the switch to the production path
and from the low-skill regime at project inception. The most recent push lifted ST
$60.0\to 65.0$ ($+5.0$pp) and HV $54.1\to 59.7$ ($+5.6$pp).

\section{What drove the recent gains}

Crucially, \emph{none} of the recent push was a reweighting of an already-priced feature. Every
increment was new data or new structure, exactly as the method predicts:
\begin{itemize}
  \item \textbf{Barrier trials} (new data): the only on-track form a debutant has; ST $+1.9$pp, with
        a full four-year panel at $94.8\%$ debut coverage.
  \item \textbf{Collateral form} (new structure --- relational opponent strength): the competition
        axis cracked by the Massey velocity rating where scalar Elo had failed; HV $+2.3$pp, ECE
        improving.
  \item \textbf{Rating peak-gap} (new structure --- a career-landmark process feature): HV $+2.4$pp,
        non-overlapping with collateral.
  \item \textbf{Trajectory fade and positional fade, recoupled to $\Wp$} (right knob): the
        reserve phenomena moved off the $\vc$ proxy onto the distance-aware reserve knob; HV
        $+0.3$pp, landing in the sprints as the $\Wp/L$ signature predicts.
  \item \textbf{The coherent trial fusion} (right structure --- own precision): the trial folded into
        the $\vc$ posterior at its own precision, fading with race count; ST $+0.4$pp overall,
        $+6$pp on debut races specifically.
  \item \textbf{Deepened collateral} (Jacobi multi-sweep + recency + trip-band): ST $+1.3$pp
        super-additive with the fusion --- the increment that broke ST through $65\%$.
  \item \textbf{Draw $\to \Ddiff$} (the third knob, first use): the trip-variance term, ST $+0.4$pp,
        ST-only as the extreme-value field-size argument predicts.
  \item \textbf{Trainer hot-streak} (intent / new structure): recency-weighted stable momentum.
\end{itemize}

\section{The scorecard: where the model sits against the theory}

Mapping the production model onto the dimensional theory of Chapter~\ref{ch:theory}, the two
highest-leverage, currently-weakest cells are \textbf{opponent-adjusted ability} (now partly
addressed by collateral form, but only its first relational moment) and \textbf{environment} ---
wind, track bias, and course bias remain entirely unused. The physics cells (kinematics, bend
geometry, surface diffusion) are in the model; the interaction cell is mean-field only; the
stochasticity cell has the diffusion noise but not the explicit start transient or in-running
incidents.

\section{The forward queue}

The ranked open levers, all framed as new data or new structure:
\begin{enumerate}[label=\textbf{Q\arabic*.},leftmargin=*]
  \item \textbf{Extend the fusion to training gallops} (queue \#30) --- a third precision channel
        $\pi_{\text{gallop}}=\rho_g\pi_{\text{race}}$ in the $\vc$ posterior; data landed
        ($30\text{k}+$ Banister sessions), needs a $\gtrsim\!1.5\text{k}$-horse backfill and a
        \code{gallop\_vc\_table}. The denser prior may convert where the sparse trial axis alone does
        not, especially at HV.
  \item \textbf{Computer-vision trajectories} (queue \#34) --- the $(x,y,t)$ feed is
        non-interceptable, so drone-replay path extraction (camera-motion homography $\to$
        rail-relative coordinates) to supply the missing lateral coordinate and unblock the
        geometry$\times$control game.
  \item \textbf{Geometry $\times$ control} (queue items around \S\ref{sec:geomgame}) --- the O3
        control-quality / ground-saved trait from existing trajectory data (the control share, no CV
        needed); draw $\times$ distance via the real bend fraction; the baked draw terms stay interim
        until geometry and control are modelled jointly.
  \item \textbf{Phase-structure levers} (queue \#31, \#32, \#33) --- the gate-break trait (start
        phase, drift + $\Ddiff$); the finishing-sectional as a high-precision $\vc$ channel
        (give each phase its own precision); the distance-law of variance, now validated on turf
        with the AWT floor identified as the only place a $D(L)$ term is warranted.
  \item \textbf{Per-horse $\Ddiff_i$ from figure variance} (queue \#28) --- a horse with erratic
        figures is genuinely less predictable; its diffusion should scale with the spread of its own
        prior figures, modulating the (constant) population baseline the distance-law study
        established. Prioritise on ST big fields and AWT per the extreme-value rule.
  \item \textbf{Environment} --- continuous soil-moisture going and wind (poll-and-log, no history);
        a deeper / directed collateral graph.
\end{enumerate}

\section{What has been ruled out (do not re-litigate as-is)}

Recorded so they are not re-tried unchanged: carried weight, rest-days (blind to training),
weight-trajectory, physiology-from-race-history (a race outcome is itself a coarse training proxy),
sustained-speed trajectory (priced), rating velocity/acceleration (restate recent form --- only the
peak-gap landmark converts), and the scalar-mean O1 draw$\times$course modulation. The standing
caveat applies: several of these are dead \emph{as a simple term} but alive as an axis with new data
or new structure (scalar Elo was ``refuted,'' yet relational collateral converts strongly), and a
parked verdict can go stale when the baseline changes --- re-check against the current production
path before trusting ``low-leverage.''

\section{Closing}

The model is a generative, physics-grounded, leak-free order-statistic model of finishing times,
calibrated to fit the market as goodness-of-fit and regularised so that every degree of freedom is a
named mechanism. It has moved from $\approx 48\%$ to $\approx 65\%/60\%$ of market by the disciplined
substitution of mechanism and new data for free parameters. The remaining gap is, on the working
hypothesis of this programme, not private information but \emph{un-modelled interaction} --- the
pace, geometry, and phase games of Chapter~\ref{ch:theory} --- and the forward queue is the staged
plan to convert each from its mean-field shadow into its structural form.

\appendix

\chapter{The baked calibration constants}
\label{app:constants}

Every value below lives in \code{hkjc.model.constants} as a mechanism-coupling gain set by the
paired fixed-seed conversion test (Chapter~\ref{ch:method}), not a free regressor weight. Per-venue
dictionaries resolve to the active value at predict time; per-surface values override the venue
value where noted.

\begin{longtable}{@{}p{0.27\textwidth}p{0.20\textwidth}p{0.45\textwidth}@{}}
\toprule
Constant & Value & Mechanism / knob \\
\midrule
\endhead
\code{FORM\_HALF\_LIFE} & $2.0$ races & Recency half-life of exp-weighted relative-finish form; the single biggest lever ($+7.8$pp). $\vc$. \\
\code{FORM\_GAIN\_BY\_VENUE} & ST $5.0$, HV $4.0$ & Form coupling; HV more weakly coupled (was over-weighted). $\vc$. \\
\code{SHRINK\_MULT} & $0.5$ & Multiplier on the sectional-$\vc$ EB residual (it was over-trusted at $1.0$). $\vc$. \\
\code{W\_PRIME\_GAIN} & $1.0$ & Amplifier on the per-horse reserve EB residual; reverted $2\to1$ once recency form absorbed the reserve signal. $\Wp$. \\
\code{JOCKEY\_GAIN} & $6.0$ & Jockey trip/ability + information proxy. $\vc$. \\
\code{DRAW\_GAIN\_BY\_VENUE} & ST $0.003$, HV $0.008$ & Bend ground-loss $\Leff = L(1+g(\mathrm{Dr}-1))$; HV bend-dominated wants $\sim\!2.7\times$ ST; zeroed on the straight chute. $\Leff$. \\
\code{DRAW\_GAIN\_AWT} & $0.008$ & Per-surface draw (dirt rail bias stronger than turf); overrides venue. $\Leff$. \\
\code{DRAW\_FS\_GAIN\_BY\_VENUE} & HV $0.25$, ST $0.0$ & Crowding $\times$ draw, clamped $\ge 0$; interim stopgap, HV-only ($+0.3$pp). $\Leff$. \\
\code{DRAW\_VAR\_GAIN\_BY\_VENUE} & ST $0.5$, HV $0.0$ & Draw $\to$ diffusion $\Ddiff$ (trip variance); ST-only by the EVT field-size argument. $\Ddiff$. \\
\code{ALLOMETRIC\_B\_BY\_VENUE} & ST $1.1$, HV $1.0$ & Allometric mass $\vc\propto m^{b-1}$; venue-invariant physiology, ST-only by leverage. $\vc$. \\
\code{POS\_FADE\_GAIN} & $0.0$ (retired) & Positional fade on $\vc$ --- retired; moved to $\Wp$. \\
\code{POSFADE\_W\_GAIN} & $-80.0$ & Front-runner reserve burn (negative); distance-aware via $\Wp/L$. $\Wp$. \\
\code{FADE\_GAIN} & $0.0$ (retired) & Fade on $\vc$ --- retired; moved to $\Wp$. \\
\code{FADE\_W\_GAIN} & $8.0$ & Late-speed-hold $\Rightarrow$ more reserve (positive). $\Wp$. \\
\code{TRIAL\_GAIN\_BY\_VENUE} & ST $1.5$, HV $1.0$ & Additive barrier-trial debut figure; superseded by the fusion when active. $\vc$. \\
\code{TRIAL\_PI\_RATIO\_BY\_VENUE} & ST $2.0$, HV $0.0$ & Coherent trial fusion precision ratio $\rho=\pi_{\text{trial}}/\pi_{\text{race}}$; no-op at $0$. $\vc$. \\
\code{COLLAT\_GAIN\_BY\_VENUE} & ST $1.0$, HV $2.5$ & Relational Massey collateral-form coupling. $\vc$. \\
\code{COLLAT\_CONSTRUCTION\_BY\_VENUE} & ST $\{3$ sweeps, hl $4$, band-par$\}$; HV $\{\}$ & Deepened collateral construction; ST $+1.3$pp. $\vc$. \\
\code{PEAKGAP\_GAIN\_BY\_VENUE} & ST $-0.015$, HV $-0.025$ & Own-peak rating deficit (process feature). $\vc$. \\
\code{TRN\_HOT\_GAIN\_BY\_VENUE} & ST $2.0$, HV $1.0$ & Trainer hot-streak / intent. $\vc$. \\
\code{SPEED\_NOISE} & $0.12$ & Velocity-noise (log-normal time); ties the diffusion. $\Ddiff$. \\
\code{awt\_diff\_mult} & $2.0$ & AWT is twice as diffusive as turf. $\Ddiff$. \\
\code{diff\_coef} & $170$ & Base diffusion coefficient ($\mathrm{m^2/s}$). $\Ddiff$. \\
\code{N\_SAMPLES} & $6000$ & Monte-Carlo samples per race. \\
\bottomrule
\end{longtable}

\chapter{The data surface}
\label{app:data}

The model is built on several leak-free panels, each acquired and persisted by \code{hkjc.data} and
keyed so the strictly-prior joins of Chapters~\ref{ch:pop}--\ref{ch:competition} are sound.

\begin{itemize}
  \item \textbf{Results (rich).} Per-runner finish position, finishing time, distance, going,
        venue, course, draw, declared weight, official rating, jockey, trainer --- the spine of the
        form, speed-figure, collateral, jockey/trainer, and geometry channels.
  \item \textbf{Sectionals.} Per-section split times (final sections $\SI{400}{m}$, first the
        remainder), the raw input to the energetics recovery (Chapter~\ref{ch:energetics}) and the
        trajectory traits.
  \item \textbf{Run-policy / trajectory.} Normalised early/late position, position gain, fade, kick
        fraction, effort --- the leak-free policy panel behind the positional-fade and the
        trajectory-shape traits, and the substrate for the (latent) control inference of
        \S\ref{sec:geomgame}.
  \item \textbf{Barrier trials.} Trial times, batch leader splits, going, distance --- the
        debut-form channel and the trial-$\vc$ fusion (Chapters~\ref{ch:features},~\ref{ch:fusion}).
        Keyed on the short stable brand.
  \item \textbf{Rating series.} The dated official-rating time series per horse, behind the
        peak-gap process feature.
  \item \textbf{Trackwork (gallops).} $30\text{k}+$ Banister training-gallop sessions, scraped and
        leak-free, the data for the gallop-fusion extension (queue \#30); not yet wired into the
        drift.
  \item \textbf{Geometry.} The venue $\times$ surface $\times$ course \code{TrackGeometry} table
        (Chapter~\ref{ch:geometry}) --- widths, home straights, and bend segments, with the bend
        fraction the draw terms read.
  \item \textbf{Odds ticks.} Open-to-close win-pool ticks for the market yardstick and closing-line
        value; the model never sees odds at fit time.
\end{itemize}

The deliberately-unmined surface --- continuous going (soil moisture), wind, track-bias / golden-rail
state, and the lateral $(x,y,t)$ trajectory --- is the environment and geometry$\times$control
frontier of Chapter~\ref{ch:theory}, and the target of the data and computer-vision items on the
forward queue.


\end{document}
